\documentclass[12pt,a4paper]{report}
\usepackage{fontspec}
\setmainfont{Liberation Serif}
\setmonofont{Liberation Mono}
\newfontfamily\cyrillicfonttt{Liberation Mono}
\usepackage{polyglossia}
\setmainlanguage{mongolian}
\usepackage{graphicx}
\usepackage[hidelinks]{hyperref}
\usepackage{geometry}
\geometry{a4paper, left=3cm, right=2cm, top=2cm, bottom=2cm}
\usepackage{titlesec}
\usepackage{tocloft}
\usepackage{setspace}
\usepackage{tikz}
\usetikzlibrary{shapes.geometric, arrows.meta, positioning, shapes.symbols, fit, shapes.multipart, calc}
\usepackage{amsmath}
\usepackage{amssymb}
\usepackage{xcolor}
\usepackage{booktabs}
\usepackage{listings}
\usepackage{pifont}

% Code listing style
\lstset{
  basicstyle=\small\ttfamily,
  breaklines=true,
  frame=single,
  numbers=left,
  numberstyle=\tiny\color{gray},
  keywordstyle=\color{blue}\bfseries,
  commentstyle=\color{gray}\itshape,
  stringstyle=\color{orange!70!black},
  showstringspaces=false,
  tabsize=4,
  captionpos=b,
}

\onehalfspacing

% Rename Lists
\renewcommand{\contentsname}{ГАРЧИГ}
\renewcommand{\listfigurename}{ЗУРГИЙН ЖАГСААЛТ}
\renewcommand{\listtablename}{ХҮСНЭГТИЙН ЖАГСААЛТ}
\renewcommand{\bibname}{НОМ ЗҮЙ}

\begin{document}

\begin{titlepage}
    % FIRST PAGE (Outer Cover)
    \thispagestyle{empty}
    \centering
    {\large\bfseries МОНГОЛ УЛСЫН ИХ СУРГУУЛЬ \\ МЭДЭЭЛЛИЙН ТЕХНОЛОГИ, ЭЛЕКТРОНИКИЙН СУРГУУЛЬ\par}
    {\large\bfseries МЭДЭЭЛЭЛ, КОМПЬЮТЕРЫН УХААНЫ ТЭНХИМ\par}
    \vspace{4cm}
    {\Large Пүрэвдоржийн Номин-Эрдэнэ\par}
    \vspace{2.5cm}
    {\Huge\bfseries Сошиал мониторингийн систем дэх агуулгын шинжилгээ, сэдэв загварчлалын шийдэл\par}
    \vspace{0.5cm}
    {\Large (Content Analysis and Topic Modeling for Social Monitoring Systems)\par}
    \vspace{1.5cm}
    {\large Прорграмм хангамж (D061302)\\
    Бакалаврын судалгааны ажил\par}
    \vfill
    Улаанбаатар \\ \the\year{} оны 3 сар
    
    \newpage
    % SECOND PAGE (Inner Cover)
    \centering
    {\large\bfseries МОНГОЛ УЛСЫН ИХ СУРГУУЛЬ \\ МЭДЭЭЛЛИЙН ТЕХНОЛОГИ, ЭЛЕКТРОНИКИЙН СУРГУУЛЬ\par}
    {\large\bfseries МЭДЭЭЛЭЛ, КОМПЬЮТЕРЫН УХААНЫ ТЭНХИМ\par}
    \vspace{6cm}
    {\Huge\bfseries Сошиал мониторингийн систем дэх агуулгын шинжилгээ, сэдэв загварчлалын шийдэл\par}
    \vspace{0.5cm}
    {\Large (Content Analysis and Topic Modeling for Social Monitoring Systems)\par}
    \vspace{1.5cm}
    {\large Прорграмм хангамж (D061302)\\
    Бакалаврын судалгааны ажил\par}
    \vspace{2.5cm}
    
    \begin{flushleft}
        \hspace{3cm} Удирдагч: \hspace{3cm} Ж.Үйтүмэн\\
        \vspace{0.5cm}
        \hspace{3cm} Гүйцэтгэсэн: \hspace{2.5cm} П.Номин-Эрдэнэ (22b1num4640)
    \end{flushleft}

    \vfill
    Улаанбаатар \\ \the\year{} оны 3 сар
\end{titlepage}

\pagenumbering{roman}
\chapter*{Зохиогчийн баталгаа}
Миний бие Пүрэвдоржийн Номин-Эрдэнэ "Сошиал мониторингийн систем дэх агуулгын шинжилгээ, сэдэв загварчлалын шийдэл" сэдэвтэй судалгааны ажлыг гүйцэтгэсэн болохыг зарлаж дараах зүйлсийг баталж байна:
\begin{itemize}
    \item Ажил нь бүхэлдээ эсвэл ихэнхдээ Монгол Улсын Их Сургуулийн зэрэг горилохоор дэвшүүлсэн болно.
    \item Энэ ажлын аль нэг хэсгийг эсвэл бүхлээр нь ямар нэг их, дээд сургуулийн зэрэг горилохоор оруулж байгаагүй.
    \item Бусдын хийсэн ажлаас хуулбарлаагүй, ашигласан бол ишлэл, зүүлт хийсэн.
    \item Ажлыг би өөрөө (хамтарч) хийсэн ба миний хийсэн ажил, үзүүлсэн дэмжлэгийг дипломын ажилд тодорхой тусгасан.
    \item Ажилд тусалсан бүх эх сурвалжид талархаж байна.
\end{itemize}
\vspace{1cm}
Гарын үсэг: \rule{4cm}{0.4pt} \\
\vspace{0.5cm}
Огноо: \rule{4cm}{0.4pt}

\newpage
\tableofcontents
\newpage
\listoffigures
\newpage
\listoftables

\clearpage
\pagenumbering{arabic}

% ================================================================
% ҮЕЧИЛСЭН ТӨЛӨВЛӨГӨӨНИЙ ХУУДАс
% ================================================================
\chapter*{Үечилсэн Төлөвлөгөө}
\addcontentsline{toc}{chapter}{ҮЕЧИЛСЭН ТӨЛӨВЛӨГӨӨ}
\vspace{0.5cm}

\begin{table}[h!]
\centering
\renewcommand{\arraystretch}{1.4}
\begin{tabular}{|c|c|p{3.5cm}|p{6cm}|c|}
\hline
\textbf{№} & \textbf{Долоо хоног} & \textbf{Үе шат} & \textbf{Хийх ажил} & \textbf{Төлөв} \\ \hline
1 & 1–4 & Судалгаа шинжилгээ & Сэдвийн хүрээнд онолын судалгаа хийх, Sumbee.mn системийн судалгаа хийх, шийдлийн аргуудыг харьцуулах & {\color{green}\checkmark} \\ \hline
2 & 4–7 & Системийн шинжилгээ, зохиомж & Шаардлага тодорхойлох, системийн архитектур болон өгөгдлийн сангийн зохиомж гаргах & {\color{green}\checkmark} \\ \hline
3 & 8–9 & Системийн хөгжүүлэлт & Data Science модулийн хөгжүүлэлт (NER, Sentiment, Topic Modeling); FastAPI backend болон Next.js frontend хөгжүүлэх & {\color{blue}Хийгдэж буй} \\ \hline
4 & 9–11 & Нэгтгэлт ба туршилт & Системд интеграци хийж нэмэлт feature хөгжүүлэх, бодит өгөгдөл дээр туршиж сайжруулах & Хүлээгдэж буй \\ \hline
5 & 12 & Урьдчилсан хамгаалалт & Явцын тайлан боловсруулалт, урьдчилсан хамгаалалт & Хүлээгдэж буй \\ \hline
6 & 13 & Засаж залруулах & Шүүмжийн дагуу алдаа засах & Хүлээгдэж буй \\ \hline
7 & 14 & Шүүмж & Шүүмжлэгчээс санал авах & Хүлээгдэж буй \\ \hline
8 & 15 & Жинхэнэ хамгаалалт & Эцсийн тайлан боловсруулалт, жинхэнэ хамгаалалт & Хүлээгдэж буй \\ \hline
\end{tabular}
\caption{Судалгааны ажлын 15 долоо хоногийн үечилсэн төлөвлөгөө}
\end{table}

\noindent\textbf{Сэдвийн агуулга:} Энэхүү бакалаврын судалгааны ажил нь сошиал мониторингийн системд цугларсан их хэмжээний текстэн өгөгдөлд агуулгын шинжилгээ хийж, тухайн өгөгдлөөс сэдэв, түлхүүр ойлголт болон шинээр үүсч буй онцлох сэдвүүдийг автоматаар илрүүлэх сэдэв загварчлалд суурилсан шийдлийг боловсруулж, бодит системд хэрэгжүүлэхэд чиглэнэ.



% ================================================================
% НЭР ТОМЬЁОНЫ ТАЙЛБАР
% ================================================================
\newpage
\chapter*{Нэр томьёоны тайлбар}
\addcontentsline{toc}{chapter}{НЭР ТОМЬЁОНЫ ТАЙЛБАР}
\begin{itemize}
    \item \textbf{Topic Modeling}: Сэдэв загварчлал — их хэмжээний текстэн өгөгдлөөс далд сэдвүүдийг автоматаар илрүүлэх аргачлал.
    \item \textbf{NLP}: Эх хэлний боловсруулалт (Natural Language Processing)
    \item \textbf{NER}: Нэрлэсэн объект таних (Named Entity Recognition) — текст дотроос хүн, байгуулллага, байршлын нэрыг ялгаж таних.
    \item \textbf{Embedding}: Текстийг олон хэмжээст тоон векторт хувиргах арга.
    \item \textbf{Sentiment Analysis}: Сэтгэл хөдлөлийн шинжилгээ — текстийн эерэг/сөрөг/саармаг чиг хандлагыг тодорхойлох.
    \item \textbf{BERTopic}: BERT embedding болон HDBSCAN clustering ашигладаг орчин үеийн сэдэв загварчлалын арга.
    \item \textbf{Hexagonal Architecture}: Портс ба Адаптер архитектур — системийн цөм логикийг гадны хэрэгслээс тусгаарлах зохион байгуулалт.
\end{itemize}

\newpage
\chapter*{Товчилсон үгс}
\begin{itemize}
    \item \textbf{LLM}: Large Language Model (Том хэлний загвар)
    \item \textbf{NER}: Named Entity Recognition (Нэрлэсэн объект таних)
    \item \textbf{API}: Application Programming Interface
    \item \textbf{DB}: Database (Өгөгдлийн сан)
    \item \textbf{BERT}: Bidirectional Encoder Representations from Transformers
    \item \textbf{UMAP}: Uniform Manifold Approximation and Projection
    \item \textbf{HDBSCAN}: Hierarchical Density-Based Spatial Clustering of Applications with Noise
    \item \textbf{XLMR}: Cross-lingual Language Model RoBERTa
    \item \textbf{GPU}: Graphics Processing Unit (Видео карт)
    \item \textbf{CORS}: Cross-Origin Resource Sharing
\end{itemize}

% ================================================================
% УДИРТГАЛ
% ================================================================
\newpage
\chapter*{Удиртгал}
\addcontentsline{toc}{chapter}{УДИРТГАЛ}

Орчин үед байгууллага, компаниуд хэрэглэгчдийнхээ санал сэтгэгдлийг сонсох, тренд болж
буй сэдвүүдийг илрүүлэх зорилгоор сошиал мониторингийн системийг өргөнөөр ашиглах
болсон. Одоогийн байдлаар Sumbee.mn зэрэг системүүд нь ихэвчлэн түлхүүр үг дээр тулгуурласан
тоон болон эерэг/сөрөг (sentiment) анализ хийдэг бөгөөд агуулгын түвшинд задлан шинжлэх,
илүү бодитой, бүтэцчилсэн дүгнэлт гаргах боломж дутмаг байна. Мөн маш их хэмжээний
өгөгдөл дээр LLM шууд ашиглах нь тооцооллын зардал маш их шаардах ба удаан (latency
өндөр) байдаг. Иймд Монгол хэл, тэр дундаа сошиал өгөгдөл дээр үр дүнтэй ажиллаж
чадах NLP (Байгалийн хэлний боловсруулалт)-д суурилсан шийдэл боловсруулах шаардлага
тулгарч байна.

Тухайлбал Нэрлэсэн объект таних (NER), Сэдэв загварчлал (Topic Modeling) болон Сэтгэл хөдлөлийн шинжилгээ (Sentiment Analysis)-ийн нэгдсэн ажлын урсгал боловсруулах шаардлага тулгарч байна.

\section*{0.1 Судалгааны үндэс ба асуудлын тодорхойлолт}

Sumbee.mn систем нь одоогоор дараах асуудлуудтай тулгарч байна:

\begin{enumerate}
    \item \textbf{Гүнзгий агуулгын шинжилгээ дутмаг:} Одоогийн систем нь түлхүүр үгт тулгуурладаг тул тухайн нийтлэл \textit{юуны тухай} яриад байгааг (сэдэв) автоматаар ялгаж чадахгүй. Жишээлбэл, \textit{"Монгол улс Оростой хэлэлцээ хийв"} гэсэн мэдээний голлох сэдэв, оролцогч талуудыг салгаж гаргах боломж одоогоор хязгаарлагдмал.
    
    \item \textbf{LLM-ийн өртөг ба хурдны асуудал:} Их хэмжээний өгөгдөл (өдөрт олон мянган нийтлэл) дээр GPT-4 зэрэг LLM API ашиглах нь 100,000 текст дээр ойролцоогоор \$200--\$300 зарцуулдаг бөгөөд хариу авах хугацаа 1--3 секунд болдог. Энэ нь бодит цагийн мониторингийн системд тохиромжгүй.
    
    \item \textbf{Монгол хэлний NLP дэмжлэг дутмаг:} Монгол хэл нь агглютинатив шинж чанартай бөгөөд нэг үгийн олон хэлбэр, галиг бичлэгийн олон хувилбар, ярианы хэллэгийн онцлог зэргээс шалтгаалан уламжлалт NLP хэрэгслүүд бага нарийвчлалтайгаар ажилладаг.
\end{enumerate}

\noindent Эдгээр асуудлыг шийдвэрлэхийн тулд шийдлийн гурван аргыг харьцуулан судалсан:

\begin{itemize}
    \item \textbf{Өгөгдлийн санд мод бүтэцтэй мэдлэгийн сан (Knowledge Tree DB):} Хамгийн их давтагдаж буй токенуудыг ялган гараар ангилал үүсгэх арга. GPU шаардахгүй, хурдан query хийх боломжтой боловч гараар ангилах их хугацаа орж, шинэ үг гарвал гараар нэмэх шаардлагатай.
    \item \textbf{Pretrained BERT загварыг fine-tune хийх:} BERT зэрэг pretrained загваруудыг монгол хэлний сошиал өгөгдөл дээр fine-tune хийж ашиглах арга. GPU дэмжлэгтэй T4 дээр 20--60 ms latency, 200--500 өгүүлбэр/секунд гүйцэтгэлтэй.
    \item \textbf{GPT / LLM API:} GPT-4o зэрэг загварыг API-аар дуудах. Сургалт шаардалгүй, маш flexible боловч хариу авах хугацаа 1--3 секунд, зардал өндөр, offline ашиглах боломжгүй.
\end{itemize}

\begin{table}[h!]
\centering
\renewcommand{\arraystretch}{1.4}
\begin{tabular}{|l|c|l|l|}
\hline
\textbf{Арга} & \textbf{GPU} & \textbf{Latency} & \textbf{Зардал} \\ \hline
Knowledge DB & Шаардалгүй & 1–5 ms & Маш бага \\ \hline
Fine-tuned BERT & Заавал биш & 20–100 ms & Бага \\ \hline
LLM API (GPT-4o) & Шаардалгүй & 1–3 sec & Өндөр (\$200–\$300/100K текст) \\ \hline
\end{tabular}
\caption{Инференсийн аргуудын харьцуулалт}
\end{table}

\noindent Судалгааны дүнгээс харахад \textbf{Pretrained BERT загвар} нь хурд, нарийвчлал, зардлын харьцааны хувьд хамгийн тохиромжтой шийдэл бөгөөд Knowledge DB-тэй хослуулан ашиглах нь бодит production pipeline-д хамгийн оновчтой гэж тогтоов.

\section*{0.2 Судалгааны зорилго ба зорилт}

\textbf{Зорилго:} Sumbee.mn сошиал мониторингийн системд Нэрлэсэн объект таних (NER), Сэдэв загварчлал (Topic Modeling), Сэтгэл хөдлөлийн шинжилгээ (Sentiment Analysis) болон Сүлжээний шинжилгээ (Network Analysis)-ийн нэгдсэн NLP pipeline боловсруулж, бодит системтэй интеграци хийснээр агуулгын гүнзгий шинжилгээний боломжийг нэмэгдүүлэхэд оршино.

\noindent\textbf{Зорилтууд:}
\begin{enumerate}
    \item Sumbee.mn-ийн одоогийн системийн бүтэц, хязгаарлалтыг судлаж, NLP pipeline нэвтрүүлэх боломжийг тодорхойлох.
    \item Монгол хэлний текст боловсруулахад тохирсон preprocessing аргачлал боловсруулах.
    \item \texttt{Davlan/bert-base-multilingual-cased-ner-hrl} загварт тулгуурлан NER pipeline хэрэгжүүлэх.
    \item BERTopic болон Cosine Similarity-д суурилсан Predefined Topic Classification аргыг хэрэгжүүлж харьцуулалт хийх.
    \item Хэрэглэгчид үр дүнг дүрслэн харуулах Next.js frontend dashboard хөгжүүлэх.
    \item FastAPI-д суурилсан Hexagonal Architecture ашиглан модульчлагдсан, дахин ашиглагдах боломжтой backend систем барих.
\end{enumerate}

\section*{0.3 Судалгааны арга зүй, хүрээ ба хязгаарлалт}

\textbf{Судалгааны арга зүй:}
Судалгаанд чанарын (qualitative) болон тоон (quantitative) шинжилгээний аргуудыг хослуулан ашиглав. Үндсэн арга зүй нь баримт бичгийн судалгаа, системийн архитектурын шинжилгээ, болон бодит өгөгдөл дээрх туршилт дээр суурилсан.

Баримт бичгийн судалгааны хүрээнд NLP-ийн суурь онол, BERTopic, BERT, XLMR загваруудын дэлгэрэнгүй техникийн баримт бичиг болон ижил системүүдийн (Brandwatch, Meltwater, Sprout Social) судалгаа хийсэн.

Хэрэгжүүлэлтийн хүрээнд Sumbee.mn-ийн бодит сошиал өгөгдөл (Facebook post, мэдээ) дээр NER, Topic Modeling, Sentiment Analysis хийж, хэрэглэгчдэд дүрслэн харуулах бүрэн end-to-end систем бүтээсэн.

\textbf{Судалгааны хүрээ:}
Судалгааны ажлын хүрээнд Монгол хэлний сошиал өгөгдлийн боловсруулалт болон Sumbee.mn системтэй интеграцид анхаарлаа хандуулав. Судалгаанд ашигласан технологиуд: Python, FastAPI, Next.js, Hugging Face Transformers, BERTopic, NetworkX, SQLite.

\textbf{Судалгааны хязгаарлалт:}
\begin{itemize}
    \item Судалгаа нь Монгол хэлний сошиал өгөгдөлд төвлөрсөн бөгөөд бусад хэлний текстэд нарийвчилсан туршилт хийгдээгүй.
    \item NER загварыг шинээр сургахгүйгээр multilingual pretrained загвар ашигласан тул Монгол хэлний NER-ийн нарийвчлал хязгаарлагдмал байж болзошгүй.
    \item Бодит production орчинд бүрэн хэрэгжүүлэлт хийх нь судалгааны ажлын дараагийн шат болно.
\end{itemize}

\chapter{Онолын судалгаа}

% ================================================================
\section{Сошиал мониторингийн системийн онолын үндэс}
% ================================================================

\subsection{Сошиал мониторингийн тухай ойлголт}

Сошиал мониторинг (Social Media Monitoring, SMM) гэдэг нь олон нийтийн сүлжээ, мэдээний сайт, блог болон бусад цахим эх сурвалжаас өгөгдөл цуглуулж, тэдгээрийг дүн шинжилгээнд оруулан бодит цагийн мэдээлэл гаргаж авах автоматжуулсан үйл явц юм~\cite{kaplan2010social}. 
Орчин үеийн сошиал мониторингийн систем нь дараах үндсэн функцуудыг гүйцэтгэдэг:
\begin{enumerate}
    \item \textbf{Өгөгдөл цуглуулалт (Data Collection):} API, web scraping болон RSS зэрэг технологи ашиглан олон эх сурвалжаас бодит цагт өгөгдөл татах.
    \item \textbf{Текст боловсруулалт (Text Processing):} Цугласан өгөгдлийг цэвэрлэж, стандартчилж, шинжилгээнд бэлтгэх.
    \item \textbf{Агуулгын шинжилгээ (Content Analysis):} NLP технологи ашиглан текстийн сэдэв, сэтгэл хөдлөл, оролцогч талуудыг тодорхойлох.
    \item \textbf{Тайлан ба Дүрслэл  (Visualization \& Reporting):} Үр дүнг хэрэглэгчид ойлгомжтой хэлбэрт оруулан харуулах.
\end{enumerate}

\subsection{Их өгөгдлийн боловсруулалтын сорилтууд}

Интернэтийн хэрэглэгчдийн тоо 2024 оны байдлаар 5.4 тэрбумд хүрч, өдөрт 500 сая tweet, 350 сая Facebook пост зэрэг хэмжээлшгүй их өгөгдөл үүсч байна~\cite{statista2024}. Энэхүү \textit{``Их өгөгдөл''} (Big Data)-ийн боловсруулалтад дараах үндсэн сорилтууд тулгардаг:

\begin{itemize}
    \item \textbf{Хурд (Velocity):} Өгөгдлийн урсгал маш их тул бодит цагийн боловсруулалт шаардагддаг.
    \item \textbf{Хэмжээ (Volume):} Хэдэн тэрбум текст бичвэрийг хадгалж, задлан шинжлэх шаардлагатай.
    \item \textbf{Олон янз байдал (Variety):} Бүтэцлэгдсэн (structured), хагас бүтэцтэй болон бүтэцлэгдээгүй (unstructured) өгөгдлийг нэгэн зэрэг боловсруулах.
    \item \textbf{Үнэн зөв байдал (Veracity):} Сошиал өгөгдөл дотор spam, давтагдсан контент их байдаг.
\end{itemize}

Эдгээр сорилтуудыг шийдвэрлэхэд NLP технологи, тэр дундаа deep learning-д суурилсан аргуудын ач холбогдол улам бүр нэмэгдэж байна.

% ================================================================
\section{Судлагдсан байдал}
% ================================================================

\subsection{Дэлхийн зах зээлийн судалгаа: Ижил чиглэлийн системүүд}

Сошиал мониторингийн чиглэлд олон улсад хэд хэдэн том систем бий болсон бөгөөд эдгээрийн архитектур, боловсруулалтын арга нь судалгааны чухал суурь болж байна.

\textbf{Brandwatch} нь зах зээлийн тэргүүлэгч сошиал мониторингийн платформ бөгөөд Apache Kafka-д суурилсан бодит цагийн stream processing pipeline ашигладаг. Их хэмжээний текстэд Latent Semantic Analysis (LSA) болон supervised machine learning загварыг хослуулан сэдэв ангилал хийдэг~\cite{brandwatch2023}. Гэвч Brandwatch нь Монгол хэлийг дэмжихгүй, enterprise зардал өндөр байдаг.

\textbf{Meltwater} нь Elasticsearch-д суурилсан real-time индексжүүлэлт ашиглан 270 гаруй хэл дэх мэдээг хянадаг. Sentiment analysis-д машин сургалтын хоёртын ангилалч (binary classifier) хэрэглэж байгаа бол сэдэв илрүүлэлтэд LDA загвараар тодорхойлсон category taxonomy ашигладаг~\cite{meltwater2022}.

\textbf{Sprout Social} нь NLP pipeline-д spaCy болон Transformers-д суурилсан multi-label classification ашиглан нийтлэлийг олон категорит нэгэн зэрэг ангилдаг.

Эдгээр системүүдийн харьцуулалтаас харахад (хүснэгт~\ref{tab:systems_comparison}), Монгол хэлний дэмжлэг болон low-latency inference нь одоогоор шийдэгдээгүй сорилт хэвээр байна.

\begin{table}[h!]
\centering
\renewcommand{\arraystretch}{1.4}
\begin{tabular}{|l|c|c|c|c|}
\hline
\textbf{Систем} & \textbf{Хэлний дэмжлэг} & \textbf{NLP арга} & \textbf{Монгол хэл} & \textbf{Нээлттэй эх} \\ \hline
Brandwatch & 25+ хэл & LSA + ML & Үгүй & Үгүй \\ \hline
Meltwater & 270+ хэл & LDA + Rule & Хязгаарлагдмал & Үгүй \\ \hline
Sprout Social & 30+ хэл & Transformer & Үгүй & Үгүй \\ \hline
Sumbee.mn (одоогийн) & Монгол & Keyword & Тийм & Үгүй \\ \hline
Энэхүү судалгаа & Монгол & BERT + BERTopic & Тийм & Тийм \\ \hline
\end{tabular}
\caption{Сошиал мониторингийн системүүдийн харьцуулалт}
\label{tab:systems_comparison}
\end{table}

\subsection{Монгол хэлний NLP-ийн судалгааны байдал}

Монгол хэлний NLP нь дэлхийн хэмжээнд харьцангуй шинэ судалгааны чиглэл бөгөөд хэд хэдэн онцлог сорилттой тулгардаг.

\textbf{Морфологийн онцлог:} Монгол хэл нь \textit{агглютинатив} (agglutinative) бүтэцтэй бөгөөд нэг үндсэн үгнээс олон зуун хэлбэр үүсэх боломжтой. Жишээлбэл, ``байгуулагдаагүйгийнх нь'' гэсэн нэг үг нь салангид морф хэсгүүдийн нийлбэрээс бүрддэг. Энэ онцлог нь уламжлалт tokenization аргуудыг үр дүнгүй болгодог~\cite{mongolnlp2020}.

\textbf{Галиг бичлэгийн олон хувилбар:} Сошиал орчинд Монгол хэрэглэгчид латин үсгийн галиглыг өргөнөөр ашигладаг. ``Монгол'', ``Mongol'', ``mongol'', ``мongol'' зэрэг олон хэлбэрийг нэг утгаар таних шаардлагатай байдаг.

\textbf{Академик судалгааны байдал:} Mongolian NLP Corpus~\cite{mongolcorpus2019} нь 1.2 сая өгүүлбэртэй корпус бүрдүүлсэн бол Ж.Пүрэв болон Ч.Одбаяр нар~\cite{purev2008} Монгол хэлний morph analysis-д суурилсан POS tagging системийг боловсруулсан. Харин сошиал өгөгдөлд зориулсан topic modeling, NER-ийн судалгаа маш хязгаарлагдмал хэвээр байна.

\textbf{Multilingual загваруудын хувь нэмэр:} XLM-RoBERTa~\cite{conneau2020xlmr}, mBERT~\cite{devlin2019bert} зэрэг multilingual pretrained загварууд Монгол хэлний downstream task-д ашиглах боломжтой болсон нь том ахиц юм. \texttt{Davlan/bert-base-multilingual-cased-ner-hrl} загвар нь 10 гаруй хэлний NER task-д fine-tune хийгдсэн бөгөөд Монгол хэлний entity recognition-д харьцангуй сайн үр дүн үзүүлж байна.

\subsection{Topic Modeling-ийн судалгааны хөгжил}

Topic modeling-ийн чиглэлд Blei болон бусад эрдэмтдийн 2003 онд санал болгосон LDA (Latent Dirichlet Allocation)~\cite{blei2003lda} нь суурь алгоритм болж тогтсон. Дараагийн томоохон ахиц нь 2022 онд Grootendorst-ын санал болгосон BERTopic~\cite{grootendorst2022bertopic} бөгөөд энэ арга нь transformer embedding-ийг dimensionality reduction болон density-based clustering-тэй хослуулан LDA-аас хамаагүй нарийвчлалтай сэдэв илрүүлэх боломжийг нээсэн.

% ================================================================
\section{Байгалийн хэлний боловсруулалт (NLP)-ийн онолын суурь}
% ================================================================

\subsection{Transformer архитектур ба BERT загвар}

2017 онд Google Brain-ийн эрдэмтэд ``Attention Is All You Need''~\cite{vaswani2017attention} өгүүлэлдээ \textit{Transformer} архитектурыг танилцуулсан нь NLP-ийн хөгжилд маш том хувьсгал болсон. Transformer нь RNN-ийн дараалалт боловсруулалтын оронд \textit{self-attention} механизмыг ашиглан өгүүлбэрийн бүх токенуудын хоорондын хамаарлыг нэгэн зэрэг тооцдог. Энэ нь:
\begin{itemize}
    \item Параллель тооцоолол хийх боломж олгодог (RNN-аас хурдан),
    \item Урт хамаарлыг (long-range dependency) зөв барьдаг,
    \item Хэлний контекстийг илүү нарийн дүрсэлдэг.
\end{itemize}

2019 онд Google-ийн судлаач Devlin болон бусад~\cite{devlin2019bert} эрдэмтэд BERT (Bidirectional Encoder Representations from Transformers) загварыг танилцуулсан. BERT нь \textit{masked language modeling} (MLM) болон \textit{next sentence prediction} (NSP) хоёр даалгаврыг нэгэн зэрэг сургаснаар текстийн хоёр талт (bidirectional) контекстийг авч чаддаг болсон нь уг загварыг NLP-ийн суурь загвар болгосон.

Загварын томьёолол дараах байдалтай:
\begin{equation}
    \text{Attention}(Q, K, V) = \text{softmax}\left(\frac{QK^T}{\sqrt{d_k}}\right)V
\end{equation}
Энд $Q$ (Query), $K$ (Key), $V$ (Value) нь загварын дотоод матриц, $d_k$ нь Key-ийн хэмжээс юм.

\subsection{Embedding ба Текстийн олон хэмжээст дүрслэл}

Текстийн embedding нь үг, өгүүлбэр болон баримт бичгийг тасралтгүй тоон векторт хөрвүүлэх арга юм. Уламжлалт bag-of-words аргаас ялгаатай нь embedding нь семантик утгыг хадгалдаг: утга ойролцоо үгс тоон орон зайд ойр байрладаг.

\textbf{Word2Vec}~\cite{mikolov2013word2vec} нь 2013 онд Mikolov-ын санал болгосон хоёр давхаргат нейрон сүлжээнд суурилсан word embedding арга бөгөөд Skip-gram болон CBOW гэсэн хоёр хувилбартай.

\textbf{Sentence-BERT (SBERT)}~\cite{reimers2019sbert} нь BERT загварт \textit{siamese network} бүтэц нэмж, өгүүлбэрийн оновчтой embedding үүсгэдэг. SBERT нь cosine similarity-г шууд тооцоолоход оновчтой бөгөөд:
\begin{equation}
    \text{sim}(u, v) = \frac{u \cdot v}{\|u\| \cdot \|v\|}
\end{equation}
гэсэн томьёогоор хоёр текстийн ойролцоо байдлыг хэмждэг. Энд $u, v \in \mathbb{R}^d$ нь текстийн embedding вектор юм.

Энэхүү судалгаанд \texttt{sentence-transformers/paraphrase-multilingual-mpnet-base-v2} загвар ашиглан Монгол текстийн semantic embedding үүсгэн topic classification болон similarity analysis-д ашигласан.

\subsection{Сэдэв загварчлал (Topic Modeling)}

\subsubsection{LDA (Latent Dirichlet Allocation)}

LDA~\cite{blei2003lda} нь Bayesian магадлалын онолд суурилсан generative загвар бөгөөд $K$ тооны далд сэдэв (latent topic) байна гэж үзэж, баримт бичиг бүр сэдвүүдийн холимог, сэдэв бүр үгсийн тархалт гэж үздэг. Загварыг тодорхойлох үндсэн магадлалын томьёолол:
\begin{equation}
    p(\mathbf{w}, \mathbf{z}, \theta, \phi | \alpha, \beta) = \prod_{k=1}^{K} p(\phi_k | \beta) \prod_{d=1}^{D} p(\theta_d | \alpha) \prod_{n=1}^{N_d} p(z_{dn} | \theta_d) p(w_{dn} | \phi_{z_{dn}})
\end{equation}
Энд $\theta_d$ нь баримт $d$-ийн сэдвийн тархалт, $\phi_k$ нь $k$-р сэдвийн үгсийн тархалт, $\alpha, \beta$ нь Dirichlet prior hyperparameter юм.

LDA-ийн давуу тал нь тайлбарлагдах боломжтой, бусад аргачлалуудтай харьцуулахад тооцоолол бага хийгддэг. Гэхдээ pre-defined topic count $K$ шаардагддаг, embedding-ийн семантик ойролцоо байдлыг тооцдоггүй, богино текстэд (tweet, пост) муу ажилладаг сул талтай.

\subsubsection{BERTopic}

BERTopic~\cite{grootendorst2022bertopic} нь орчин үеийн сэдэв загварчлалын арга бөгөөд дараах дөрвөн үе шатаас бүрдэнэ:

\begin{enumerate}
    \item \textbf{Embedding үүсгэх:} Sentence-BERT загвараар текст бүрийг $\mathbb{R}^{768}$ хэмжээст вектор болгон хөрвүүлэх.
    \item \textbf{Хэмжээс бууруулалт:} UMAP (Uniform Manifold Approximation and Projection) ашиглан embedding-ийг $\mathbb{R}^{5}$ болгон бууруулах.
    \item \textbf{Кластерлалт:} HDBSCAN (Hierarchical Density-Based Spatial Clustering) ашиглан нягт бүлэглэлийг илрүүлэх, дуу чимээ (noise) гэж тодорхойлсон цэгүүдийг \texttt{-1} ангилалд хамааруулах.
    \item \textbf{Сэдвийн дүрслэл:} c-TF-IDF (class-based TF-IDF) ашиглан сэдэв бүрийн хамгийн онцлог үгсийг гаргах.
\end{enumerate}

BERTopic-ийн c-TF-IDF томьёо:
\begin{equation}
    W_{t,c} = tf_{t,c} \cdot \log\left(1 + \frac{A}{tf_t}\right)
\end{equation}
Энд $tf_{t,c}$ нь $c$ классийн $t$ үгийн давтамж, $A$ нь нийт баримт бичгийн дундаж тоо, $tf_t$ нь нийт класс дахь $t$ үгийн давтамж юм.

\subsubsection{Predefined Topic Classification}

Урьдчилан тодорхойлсон сэдвийн ангилалт нь сэдэв бүрт жишиг embedding үүсгэж, шинэ текстийн embedding-тэй cosine similarity харьцуулан хамгийн төстэй сэдэвт ангилдаг. Энэ арга нь:
\begin{itemize}
    \item inference хурд маш өндөр (1--5 ms/текст),
    \item сэдвийн тоо ба агуулгыг бизнесийн хэрэгцээнд тохируулах боломжтой,
    \item шинэ сэдэв нэмэхэд дахин сургалт шаардахгүй.
\end{itemize}
Харин шинэ, урьдчилан тодорхойлогдоогүй сэдвүүдийг илрүүлж чаддаггүй нь уг аргын гол хязгаарлалт юм.

\subsection{Нэрлэсэн объект таних (Named Entity Recognition, NER)}

NER нь текст дотроос тодорхой ангиллын нэрлэсэн объектуудыг (named entities) автоматаар таних sequence labeling даалгавар юм. Уламжлалт NER нь PERSON, ORG, LOC, DATE, GPE зэрэг стандарт ангилал ашигладаг.

\textbf{Sequence labeling:} NER даалгаварт BIO (Begin-Inside-Outside) tagging schema өргөнөөр хэрэглэгддэг:
\begin{itemize}
    \item \texttt{B-PER}: хүний нэрийн эхний токен
    \item \texttt{I-PER}: хүний нэрийн дотоод токен
    \item \texttt{O}: entity биш токен
\end{itemize}

BERT-д суурилсан NER нь токен бүрт softmax ангилалч ашиглан BIO шошго оноодог:
\begin{equation}
    p(y_i | x_i, \mathbf{h}) = \text{softmax}(W \cdot h_i + b)
\end{equation}
Энд $h_i$ нь $i$-р токений BERT-ийн hidden state, $W$ ба $b$ нь сургагдах параметр юм.

Энэхүү судалгаанд \texttt{Davlan/bert-base-multilingual-cased-ner-hrl} загварыг ашигласан бөгөөд энэ нь 10 хэлний CoNLL-2003, WikiAnn зэрэг стандарт NER датасет дээр fine-tune хийгдсэн multilingual загвар юм.

\subsection{Сэтгэл хөдлөлийн шинжилгээ (Sentiment Analysis)}

Sentiment analysis нь текстийн субьектив санал бодлын чиглэлийг (positive, negative, neutral) автоматаар тодорхойлох даалгавар юм. Уг чиглэлд хоёр үндсэн хандлага байдаг:

\begin{enumerate}
    \item \textbf{Lexicon-based арга:} Урьдчилан бүрдүүлсэн sentiment dictionary ашиглан үг бүрт оноо өгч, нийтлэг байдлыг тодорхойлдог. Монгол хэлний хувьд sentiment lexicon дутмаг байдаг нь энэ аргын гол сул тал.
    \item \textbf{Machine Learning-д суурилсан арга:} Сургагдсан датасет дээр ангилалч загвар суралцуулдаг. BERT-д суурилсан загварууд нь уг аргын тэргүүлэгч шийдэл болжээ.
\end{enumerate}

Sentiment ангиллын алдааг хэмжихэд $F_1$ score ашиглагддаг:
\begin{equation}
    F_1 = \frac{2 \cdot \text{Precision} \cdot \text{Recall}}{\text{Precision} + \text{Recall}}
\end{equation}

Монгол хэлний сошиал өгөгдөлд \texttt{tabularisai/multilingual-sentiment-analysis} загвар ашигласан бөгөөд энэ нь 8 хэлний sentiment датасет дээр fine-tune хийгдсэн.

\subsection{Сүлжээний шинжилгээ (Network Analysis)}

Сүлжээний шинжилгээ нь NER-ийн үр дүнд гарсан entity-үүдийн хамт дурдагдсан байдал (co-occurrence) буюу хоорондын хамаарлыг граф хэлбэрт хөрвүүлэн судлах арга юм. Граф $G = (V, E)$ болон:
\begin{itemize}
    \item $V$: NER-ийн entity-үүд (хүмүүс, байгуулллага, байршил)
    \item $E$: Нэг баримт бичигт хамт дурдагдсан entity-үүдийн хоорондын ирмэг (edge)
    \item $w(e)$: Хамт дурдагдалтын давтамж (edge weight)
\end{itemize}

Сүлжээний шинжилгээнд дараах үндсэн хэмжигдэхүүнүүдийг тооцдог:
\begin{equation}
    C_D(v) = \frac{\deg(v)}{|V| - 1} \quad \text{(Degree Centrality)}
\end{equation}
\begin{equation}
    C_B(v) = \sum_{s \neq v \neq t} \frac{\sigma_{st}(v)}{\sigma_{st}} \quad \text{(Betweenness Centrality)}
\end{equation}
Энд $\sigma_{st}$ нь $s$-аас $t$ руу хамгийн богино замын тоо, $\sigma_{st}(v)$ нь $v$ оройгоор дайрах тийм замын тоо юм. Betweenness centrality өндөр entity нь мэдээллийн сүлжээний ``гүүр'' буюу хамгийн нөлөөтэй зангилаа болдог.

% ================================================================
\section{Sumbee.mn системийн архитектур ба одоогийн байдлын шинжилгээ}
% ================================================================

\subsection{Sumbee.mn системийн тухай}

Sumbee.mn нь Монгол улсад 2016 оноос үйл ажиллагаа явуулж буй сошиал мониторингийн платформ юм. Систем нь Facebook, Twitter/X, болон монгол дахь үйл ажиллагаа явуулж буй мэдээний сайт, блогуудаас өгөгдөл цуглуулан хэрэглэгчдэд бодит цагийн мониторинг, сентимент шинжилгээ болон тренд анализийн мэдээлэл өгдөг. Системийн үйлчлүүлэгчдэд хувь хүн, албан байгуулллага, бизнес эрхлэгчид багтдаг.

\subsection{Одоогийн системийн ерөнхий архитектур}

Sumbee.mn системийн одоогийн техник бүтэц нь зураг~\ref{fig:sumbee_arch}-д үзүүлснээр хэд хэдэн давхаргаас бүрдэнэ. 
\begin{figure}[h!]
\centering
\resizebox{\linewidth}{!}{%
\begin{tikzpicture}[
    node distance=0.8cm and 1.2cm,
    every node/.style={font=\small},
    server/.style={draw, rounded corners=3pt, fill=blue!10, minimum width=2.8cm, minimum height=1.2cm, text centered, text width=2.6cm},
    db/.style={draw, cylinder, shape border rotate=90, fill=cyan!20, minimum width=1.6cm, minimum height=1.2cm, text centered, text width=1.4cm},
    process/.style={draw, rounded corners=3pt, fill=orange!15, minimum width=3.2cm, minimum height=1.5cm, text centered, text width=3cm},
    app/.style={draw, rounded corners=3pt, fill=green!15, minimum width=2.8cm, minimum height=1.2cm, text centered, text width=2.6cm},
    cloud/.style={draw, ellipse, fill=gray!15, minimum width=2cm, minimum height=1cm, text centered, text width=1.8cm},
    box/.style={draw, dashed, rounded corners=5pt, inner sep=8pt},
    arrow/.style={->, >=stealth, thick}
]

% --- Crawlers (left) ---
\node[server] (tw) {TW Crawler\\};
\node[server, below=0.5cm of tw] (fb) {FB Crawler\\};
\node[server, below=0.5cm of fb] (trend) {YT Crawler\\};
\node[server, below=0.5cm of trend] (web) {Web Crawler\\};

\node[box, fit=(tw)(web), label=above:{\textbf{Crawler давхарга}}] (crawlerbox) {};

% --- Raw DB ---
\node[db, right=1.5cm of fb] (rawdb) {Text data};
\node[server, below=0.8cm of rawdb] (media1) {Media data};

\node[box, fit=(rawdb)(media1), label=above:{\textbf{Raw Storage}}] (rawbox) {};

% --- Processing ---
\node[process, right=1.5cm of rawdb] (proc) {Processing\\Server\\NLP Pipeline\\};

% --- Processed DBs ---
\node[db, above right=0.35cm and 1.2cm of proc] (db3p) {DB\\User specific};
\node[db, below right=0.35cm and 1.2cm of proc] (db2p) {DB\\General};

\node[box, fit=(db3p)(db2p), label=above:{\textbf{Боловсруулсан өгөгдлийн сан}}] (dbbox) {};

% --- Web Server ---
\node[server] (web33) at ($(db3p)!0.5!(db2p) + (3.4cm, 0)$) {Web Server\\(Application)};

% --- API ---
\node[cloud, below=1cm of proc] (api) {API / GPT\\(Cloud)};

% --- Arrows ---
\draw[arrow] (tw) -- (rawdb);
\draw[arrow] (fb) -- (rawdb);
\draw[arrow] (trend) -- (rawdb);
\draw[arrow] (web) -- (rawdb);
\draw[arrow] (tw) -- (media1);
\draw[arrow] (fb) -- (media1);

\draw[arrow] (rawdb) -- (proc);
\draw[arrow] (media1) -- (proc);

\draw[arrow] (proc) -- (db3p);
\draw[arrow] (proc) -- (db2p);
\draw[arrow] (proc) -- (api);
\draw[arrow] (api) -- (proc);

\draw[arrow] (db2p) -- (web33);
\draw[arrow] (db3p) -- (web33);

\end{tikzpicture}%
}
\caption{Sumbee.mn (СИТНЭТ) системийн ерөнхий архитектур}
\label{fig:sumbee_arch}
\end{figure}

\subsection{Системийн давхаргуудын шинжилгээ}

\textbf{Crawler давхарга:} Олон төрлийн crawler (Twitter/X, Facebook, YouTube, Web) зэрэгцэн (parallel) ажиллаж олон нийтийн сүлжээ болон мэдээний эх сурвалжаас өгөгдөл цуглуулдаг.

\textbf{Raw Storage давхарга:} Crawler-уудаас ирсэн текст болон медиа (зураг, видео) өгөгдлийг тусдаа хадгалалтын бүрэлдэхүүнүүдэд хадгалдаг.

\textbf{Processing Server:} Цугласан өгөгдлийг текст ангилал, шошголол, сэдэв тодорхойлолт, медиа боловсруулалт зэрэг NLP модулиудаар дамжуулан боловсруулдаг. Шаардлагатай тохиолдолд гадаад AI API-тай холбогдох боломжтой.

\textbf{Боловсруулсан өгөгдлийн сан:} Боловсруулалтын үр дүнг хэрэглэгчийн хэрэгцээнд тохирсон бүтэцтэй өгөгдлийн санд хадгалдаг.

\textbf{Application давхарга:} Хэрэглэгчийн интерфейс болон бизнесийн логикийг агуулсан вэб сервер бөгөөд боловсруулсан өгөгдлийг хэрэглэгчдэд харуулдаг.

\subsection{Одоогийн системийн хязгаарлалтууд}

Sumbee.mn-ийн одоогийн Processing Server нь текст боловсруулалтад ихэвчлэн keyword matching болон rule-based арга хэрэглэдэг. Энэ нь:

\begin{itemize}
    \item \textbf{Сэдвийн автомат тодорхойлол дутмаг:} Нийтлэл \textit{юуны тухай} яриад байгааг (семантик сэдэв) автоматаар ялгаж чаддаггүй.
    \item \textbf{Entity-ийн харилцааны шинжилгээ дутмаг:} Ямар хүн, байгуулллага ямар сэдвийн хүрээнд хамт дурдагдаж байгааг тодорхойлох чадваргүй.
    \item \textbf{LLM-ийн тооцооллын зардал:} Их хэмжээний өгөгдөл дээр GPT-4 зэрэг LLM API ашиглах нь 100,000 текстэд \$200--\$300, latency 1--3 секунд болдог тул бодит цагийн системд тохиромжгүй.
\end{itemize}

Энэхүү судалгааны NLP pipeline нь уг хязгаарлалтуудыг шийдвэрлэхэд чиглэсэн бөгөөд Processing Server-ийн шинэ модуль болон нэмэгдэх юм.

% ================================================================
\section{Аргуудын харьцуулалт ба шийдлийн сонголт}
% ================================================================

\subsection{Текст боловсруулалтын гурван аргын харьцуулалт}

Сошиал мониторингийн системд NLP pipeline нэвтрүүлэхэд дараах гурван үндсэн аргыг авч үзлээ:

\textbf{1. Мэдлэгийн санд суурилсан арга (Knowledge Base / Rule-based):}
Хамгийн их давтагдаж буй токенуудаас гараар ангилал үүсгэж, шинэ текстийг тухайн мэдлэгийн сантай харьцуулан ангилдаг. Энэ арга нь GPU шаардахгүй, query хурд маш өндөр (1--5 ms), хэрэгжүүлэхэд хялбар. Гэхдээ гараар бүрдүүлэх, шинэчлэх их хугацаа шаарддаг, шинэ үг, нэр томьёо гарахад автоматаар ажилладаггүй сул талтай.

\textbf{2. Pretrained BERT загварт суурилсан арга:}
Bert болон multilingual загваруудыг тухайн даалгаварт зохицуулан fine-tune хийж ашигладаг. T4 GPU-тэй орчинд 20--60 ms latency, 200--500 өгүүлбэр/секунд гүйцэтгэлтэй. Сургалтын датасет шаардагдах боловч нэг удаа сургасан загварыг дахин сургахгүйгээр ашиглаж болно. Хурд болон нарийвчлалын оновчтой тэнцвэрийг хангадаг.

\textbf{3. LLM API (GPT-4o зэрэг):}
Anthropic, OpenAI зэрэг компаниудын API-аар дуудаж flexible, олон даалгаварт ашиглагддаг. Сургалт шаардалгүй, шинэ хэлэнд хурдан дасан зохицдог. Гэхдээ 1--3 секундын latency, \$200--\$300/100K текстийн зардал, internet connectivity шаардлага нь бодит production системд хязгаарлалт болдог.

\begin{table}[h!]
\centering
\renewcommand{\arraystretch}{1.5}
\begin{tabular}{|l|c|c|c|}
\hline
\textbf{Хэмжүүр} & \textbf{Knowledge DB} & \textbf{Fine-tuned BERT} & \textbf{LLM API} \\ \hline
GPU шаардлага & Шаардалгүй & Заавал биш & Шаардалгүй \\ \hline
Latency & 1--5 ms & 20--100 ms & 1--3 сек \\ \hline
100K текстийн зардал & Маш бага & Бага & \$200--\$300 \\ \hline
Нарийвчлал & Дунд & Өндөр & Маш өндөр \\ \hline
Offline ажиллагаа & Тийм & Тийм & Үгүй \\ \hline
Шинэ сэдэв илрүүлэлт & Хязгаарлагдмал & Дунд & Өндөр \\ \hline
Монгол хэл & Тохируулж болно & Дунд & Дунд \\ \hline
\end{tabular}
\caption{Текст боловсруулалтын аргуудын нарийвчилсан харьцуулалт}
\label{tab:method_comparison}
\end{table}

\subsection{Шийдлийн сонголтын үндэслэл}

Хүснэгт~\ref{tab:method_comparison}-ийн харьцуулалтаас дараах дүгнэлт хийгдэв:

\begin{enumerate}
    \item \textbf{Pretrained BERT загвар} нь latency болон нарийвчлалын тэнцвэрийн хувьд хамгийн тохиромжтой шийдэл болж байна. Sumbee-ийн Processing Server-т нэмэлт GPU шаардалгүйгээр CPU орчинд ч 100--300 ms-д ажиллах боломжтой.
    \item \textbf{Knowledge Base} аргыг BERT-тэй хослуулан ашиглах нь: тодорхой, урьдчилан мэдэгдэж буй нэр томьёо, байгуулллагын нэрийг маш хурдан (1--5 ms) таних боломж олгоно.
    \item \textbf{LLM API} нь зардал болон latency-ийн асуудлаас шалтгаалан batch processing-д биш, харин тусгай дүн шинжилгээний хэрэгцээнд (жишээ нь тайлан боловсруулалт) ашиглах нь оновчтой.
\end{enumerate}

Тиймээс энэхүү судалгаанд \textbf{``Pretrained BERT + Knowledge Base''}-ийн хосолсон арга зүйг сонгосон нь зардал, хурд, нарийвчлалын гурван шаардлагыг нэгэн зэрэг хангах боломж олгодог.

% ================================================================
\section{Дүгнэлт}
% ================================================================

Онолын судалгааны бүлгийн дүгнэлтэнд:

\begin{enumerate}
    \item Сошиал мониторингийн системд их хэмжээний бүтэцлэгдээгүй текст өгөгдлийг боловсруулах NLP технологийн хэрэгцээ улам бүр нэмэгдэж байна. Монгол хэлний хувьд агглютинатив шинж, галиг бичлэгийн олон хувилбар зэрэг онцлогоос шалтгаалан NLP шийдлүүд хязгаарлагдмал байна.

    \item Академик болон дэлхийн зах зээлийн ижил системүүдийн судалгаанаас харахад BERTopic, multilingual BERT загварууд нь олон хэлний текст шинжилгээнд тэргүүлэгч арга болж тогтоогдсон. Гэвч Монгол хэлийг тусгайлан дэмждэг нарийвчилсан шийдэл дутмаг байна.

    \item Sumbee.mn системийн одоогийн архитектурын шинжилгээнд дөрвөн давхаргат (crawler, raw storage, processing, application) бүтэц тодорхойлогдсон бөгөөд Processing Server-т NLP pipeline нэмэх техникийн боломж байгааг баталгаажуулав.

    \item Аргуудын харьцуулалтаас ``Pretrained BERT + Knowledge Base'' хосолсон арга нь хурд, нарийвчлал, зардлын хувьд хамгийн тохиромжтой шийдэл болохыг тогтоов. Энэ арга нь дараах бүлгийн системийн зохиомж болон хэрэгжүүлэлтийн үндэс болно.
\end{enumerate}

\chapter{Системийн шинжилгээ ба зохиомж}

 Энэхүү судалгаа болон үр дүнг нь зөвхөн Sumbee.mn дотоод хэрэгцээнд зориулсан хаалттай хэрэгжүүлэлт биш, харин гадаад судлаачид, хөгжүүлэгчид ашиглах, туршилт хийх боломжтой \textbf{нээлттэй Public API} болгон зохиомжлох нь зүйтэй юм. Ингэснээр:

\begin{itemize}
    \item Судлаачид, өгөгдлийн шинжээчид NLP pipeline-ийг өөрийн датасет дээр туршиж, гаралтыг шинжлэх боломжтой болно.
    \item Гуравдагч тал хөгжүүлэгчид public API - ийг шууд дуудаж ажиллуулан турших хэрэглэх боломжтой болгох.
    \item Sumbee.mn-ийн Processing Server-т NLP модуль нэмэх интеграцийн ажил болон нээлттэй API-ийн ажил ижил хөгжүүлэлтийн суурь дээр явагдна.
\end{itemize}

Энэ шаардлагыг хангах зохитой архитектурын загвар бол \textbf{Hexagonal Architecture} (Ports \& Adapters) юм. Энэ загварт NLP Domain Core нь FastAPI, Next.js, Sumbee-ийн дэд бүтцийн талаар мэддэггүй, цэвэр бизнес логик хэвээр байна. Гадны технологи зөвхөн ``Adapter'' дамжуулан холбогдох тул нэг NLP цөмийг гурван хувилбарт дахин ашиглах боломж бий болдог:

\begin{itemize}
    \item Судлаачид болон өгөгдлийн шинжээчид NLP pipeline-ийг өөрсдийн датасет дээр туршин, үр дүнг нарийвчлан шинжлэх боломж бүрдэнэ.
    \item Гуравдагч талын хөгжүүлэгчид public API-ийг шууд дуудаж, турших болон ашиглах боломжтой болно.
    \item Sumbee.mn-ийн Processing Server-т NLP модуль нэвтрүүлэх интеграцийн ажил болон нээлттэй API хөгжүүлэлт нь нэг ижил хөгжүүлэлтийн суурин дээр хийгдэнэ.
\end{itemize}
Энэ бүлэгт системийн шаардлагуудыг тодорхойлж, дээрх гурван хувилбарын архитектурыг диаграммаар дүрслэн, Hexagonal зохиомж, NLP pipeline-ийн мэдээллийн урсгал, дарааллын диаграмм болон өгөгдлийн сангийн бүтцийг танилцуулна.

% ================================================================
\section{Системийн шаардлагын шинжилгээ}
% ================================================================

\subsection{Функциональ шаардлагууд}

Системийн функциональ шаардлагуудыг Sumbee.mn системд хийсэн шаардлагын шинжилгээ болон NLP pipeline-ийн техникийн боломжид тулгуурлан тодорхойллоо. 
\begin{table}[h!]
\centering
\renewcommand{\arraystretch}{1.5}
\begin{tabular}{|c|c|p{2.5cm}|p{5cm}|p{1.8cm}|c|}
\hline
\textbf{№} & \textbf{Код} & \textbf{Нэр} & \textbf{Тайлбар} & \textbf{Хэрэглэгч} & \textbf{Эрэмбэ} \\ \hline
1 & F-01 & Өгөгдөл оруулах & CSV файл эсвэл шууд текст оруулан NLP шинжилгээ эхлүүлэх. Нэг удаад олон мянган мөр хүлээн авна. & Шинжээч & Өндөр \\ \hline
2 & F-02 & Урьдчилан боловсруулалт & Монгол/галиг бичлэг, emoji, тусгай тэмдэгт, стоп үгсийг автоматаар цэвэрлэж стандартчилах. & Систем & Өндөр \\ \hline
3 & F-03 & Нэрлэсэн объект таних & Текстээс PERSON, ORG, LOC нэрлэсэн объектуудыг байрлал болон итгэлцлийн оноотойгоор буцаах. & Шинжээч & Өндөр \\ \hline
4 & F-04 & Сэтгэл хөдлөлийн шинжилгээ & Текст бүрийг Positive, Negative, Neutral ангилалд хамааруулж тархалтын оноог буцаах. & Шинжээч & Өндөр \\ \hline
5 & F-05 & Сэдэв загварчлал & BERTopic болон Predefined Classification аргаар нийтлэлүүдийг сэдэвт ангилж гол үгсийг гаргах. & Шинжээч & Өндөр \\ \hline
6 & F-06 & Сүлжээний шинжилгээ & Entity-үүдийн co-occurrence сүлжээг тооцоолж Degree болон Betweenness Centrality-г хэмжих. & Шинжээч & Дунд \\ \hline
7 & F-07 & Хяналтын самбар & Үр дүнг Pie chart, Treemap, Network graph, Timeline интерактив дүрслэлээр харуулах. & Шинжээч & Өндөр \\ \hline
8 & F-08 & Тайлан үүсгэх & Шинжилгээний нэгтгэсэн тайланг JSON/PDF хэлбэрт гаргах. & Шинжээч & Дунд \\ \hline
9 & F-09 & REST API endpoint & Гадаад систем болон Sumbee-ийн Processing Server-ээс HTTP дуудлагаар NLP үйлчилгээ авах. & Sumbee систем & Өндөр \\ \hline
10 & F-10 & Мэдлэгийн сан удирдах & Админ стоп үгс, сэдвийн ангилал, шошго нэмэх, засах, устгах. Predefined topic embedding шинэчлэх. & Админ & Дунд \\ \hline
\end{tabular}
\caption{Системийн функциональ шаардлагууд}
\label{tab:functional_req}
\end{table}

\subsection{Функциональ бус шаардлагууд}

\begin{table}[h!]
\centering
\renewcommand{\arraystretch}{1.5}
\begin{tabular}{|c|c|p{2.2cm}|p{5.5cm}|p{2.5cm}|}
\hline
\textbf{№} & \textbf{Код} & \textbf{Нэр} & \textbf{Тайлбар} & \textbf{Хэмжүүр} \\ \hline
1 & NF-01 & Архитектур & Hexagonal Architecture загварт NLP Domain Core-г FastAPI, Next.js-ээс тусгаарлах. & Coupling = 0 \\ \hline
2 & NF-02 & Гүйцэтгэл & GPU: 20--60 ms, CPU: 300--800 ms. 1000 өгүүлбэрийн batch 2--4 сек. & Latency $<$ 100 ms \\ \hline
3 & NF-03 & REST API & CORS зөв тохируулагдсан, JSON хариу, OpenAPI doc автоматаар үүсэх. & HTTP 200 + JSON \\ \hline
4 & NF-04 & Монгол хэл & Кирилл, галиг, ярианы хэллэг, emoji зөв боловсруулах preprocessing. & F1 $>$ 0.80 \\ \hline
5 & NF-05 & Өргөтгөх боломж & Шинэ NLP модуль нэмэхэд одоогийн кодыг өөрчлөхгүйгээр Plug-in хэлбэрт нэмэх. & Open/Closed зарчим \\ \hline
6 & NF-06 & Мэдлэгийн сан & Knowledge Base-ийг SQLite-д хадгалж migration, backup боломжтой байх. & SQLite 3.x \\ \hline
7 & NF-07 & UI/UX & Glassmorphic dark-mode, responsive layout, 3 секунд дотор үр дүн харуулах. & PageSpeed $>$ 80 \\ \hline
8 & NF-08 & Аюулгүй байдал & API key authentication, rate limiting, input validation, SQL injection хамгаалалт. & OWASP Top 10 \\ \hline
\end{tabular}
\caption{Системийн функциональ бус шаардлагууд}
\label{tab:nonfunctional_req}
\end{table}

\subsection{Use Case диаграмм}

Системийн хэрэглэгчийн хэрэглээний тохиолдлуудыг зураг~\ref{fig:usecase}-д үзүүлэв. Гурван актор байна: \textbf{Шинжээч} хэрэглэгч, системийн \textbf{Админ}, болон гадаад \textbf{Sumbee систем}.

\begin{figure}[h!]
\centering
\resizebox{0.95\linewidth}{!}{%
\begin{tikzpicture}[align=center,
    actor/.style={rectangle, draw, fill=blue!10, rounded corners=3pt,
        minimum width=2cm, minimum height=0.7cm, text centered, font=\small},
    usecase/.style={ellipse, draw, fill=green!10, minimum width=3.4cm,
        minimum height=0.75cm, text centered, font=\small},
    sysbox/.style={rectangle, draw=black!40, dashed, rounded corners=6pt, inner sep=12pt},
    conn/.style={thin, draw=black}
]
\node[sysbox, minimum width=9.5cm, minimum height=10cm] at (5.5,-4.5) {};
\node[font=\small\bfseries] at (5.5, 0.5) {NLP Intelligence систем};

\node[actor] (analyst) at (-0.5,-1.5) {Шинжээч};
\node[actor] (admin)   at (-0.5,-7.5) {Админ};
\node[actor] (sumbee)  at (11.5,-4.0) {Sumbee\\систем};

\node[usecase] (uc1) at (3.5,-0.8) {Текст/CSV оруулах};
\node[usecase] (uc2) at (3.5,-2.3) {NER шинжилгээ};
\node[usecase] (uc3) at (3.5,-3.8) {Sentiment шинжилгээ};
\node[usecase] (uc4) at (7.5,-0.8) {Сэдэв загварчлал};
\node[usecase] (uc5) at (7.5,-2.3) {Сүлжээний шинжилгээ};
\node[usecase] (uc6) at (7.5,-3.8) {Dashboard харах};
\node[usecase] (uc7) at (5.5,-5.4) {Тайлан гаргах};
\node[usecase] (uc8) at (3.5,-7.0) {Stop word удирдах};
\node[usecase] (uc9) at (7.5,-7.0) {Сэдвийн ангилал};
\node[usecase] (uc10) at (5.5,-8.5) {Шошго удирдах};
\node[usecase] (uc11) at (8.5,-4.8) {REST API дуудах};

\draw[conn] (analyst)--(uc1); \draw[conn] (analyst)--(uc2);
\draw[conn] (analyst)--(uc3); \draw[conn] (analyst)--(uc4);
\draw[conn] (analyst)--(uc5); \draw[conn] (analyst)--(uc6);
\draw[conn] (analyst)--(uc7);
\draw[conn] (admin)--(uc8);   \draw[conn] (admin)--(uc9);
\draw[conn] (admin)--(uc10);
\draw[conn] (sumbee)--(uc11);
\draw[conn] (uc11)--(uc2);   \draw[conn] (uc11)--(uc3);
\draw[conn] (uc11)--(uc4);
\end{tikzpicture}}
\caption{Системийн Use Case диаграмм}
\label{fig:usecase}
\end{figure}

% ================================================================
\section{Системийн гурван хэрэгжүүлэлтийн хувилбар}
% ================================================================

Hexagonal Architecture-ийн ачаар NLP Domain Core нь нэмэлт өөрчлөлтгүйгээр гурван хувилбарт адилхан ажилладаг.

\subsection{Demo Web Application архитектур}

Demo Web Application нь судалгааны үр дүнг харуулах зориулалттай бүрэн end-to-end систем бөгөөд Next.js frontend болон FastAPI backend-ийн нэгдмэл архитектуртай (зураг~\ref{fig:demo_arch}).

\begin{figure}[h!]
\centering
\resizebox{0.82\linewidth}{!}{%
\begin{tikzpicture}[
    layer/.style={rectangle, draw, rounded corners=4pt,
        minimum width=10cm, minimum height=1.1cm, text centered, font=\small\bfseries},
    sub/.style={rectangle, draw, rounded corners=3pt, fill=white,
        minimum width=2.4cm, minimum height=0.75cm, text centered, font=\footnotesize},
    arrow/.style={->, >=stealth, thick}
]
\node[layer, fill=blue!15]   (fe)   at (0, 0)    {};
\node[layer, fill=orange!15] (api)  at (0,-2.0)  {};
\node[layer, fill=green!15]  (core) at (0,-4.0)  {};
\node[layer, fill=purple!10] (kb)   at (0,-6.0)  {};

\node[font=\small\bfseries] at (-3.5, 0)   {Next.js Frontend};
\node[font=\small\bfseries] at (-3.5,-2.0) {FastAPI Backend};
\node[font=\small\bfseries] at (-3.5,-4.0) {NLP Domain Core};
\node[font=\small\bfseries] at (-3.5,-6.0) {Knowledge Layer};

\node[sub] at (-1.2, 0)   {Dashboard UI};
\node[sub] at (1.3,  0)   {Charts/Graphs};
\node[sub] at (3.8,  0)   {Network View};
\node[sub] at (-1.2,-2.0) {API Routers};
\node[sub] at (1.3, -2.0) {Pydantic Schema};
\node[sub] at (3.8, -2.0) {CORS / Auth};
\node[sub] at (-2.2,-4.0) {Preprocessor};
\node[sub] at (0.3, -4.0) {NER Engine};
\node[sub] at (2.5, -4.0) {Sentiment};
\node[sub] at (4.5, -4.0) {Topic Modeler};
\node[sub] at (-1.5,-6.0) {Network Analyzer};
\node[sub] at (1.5, -6.0) {KB Manager};
\node[sub] at (4.0, -6.0) {SQLite DB};

\draw[arrow] (0,-0.55)  -- node[right,font=\footnotesize]{REST/JSON} (0,-1.45);
\draw[arrow] (0,-2.55)  -- node[right,font=\footnotesize]{Port call}  (0,-3.45);
\draw[arrow] (0,-4.55)  -- node[right,font=\footnotesize]{SQL query} (0,-5.45);

\node[font=\footnotesize,text=gray] at (5.7, 0)   {:3000};
\node[font=\footnotesize,text=gray] at (5.7,-2.0) {:8000};
\end{tikzpicture}}
\caption{Demo Web Application-ийн давхаргат архитектур}
\label{fig:demo_arch}
\end{figure}

\subsection{Public REST API архитектур}

Public REST API хувилбар нь гадаад систем болон хөгжүүлэгчдэд NLP үйлчилгээ нийлүүлэх зориулалттай (зураг~\ref{fig:api_arch}).

\begin{figure}[h!]
\centering
\resizebox{0.88\linewidth}{!}{%
\begin{tikzpicture}[align=center,
    client/.style={rectangle, draw, fill=gray!15, rounded corners=3pt,
        minimum width=2.2cm, minimum height=0.8cm, text centered, font=\footnotesize},
    ep/.style={rectangle, draw, fill=orange!15, rounded corners=3pt,
        minimum width=3.6cm, minimum height=0.65cm, text centered, font=\footnotesize},
    mod/.style={rectangle, draw, fill=green!15, rounded corners=3pt,
        minimum width=2.4cm, minimum height=0.65cm, text centered, font=\footnotesize},
    arrow/.style={->, >=stealth, thick},
    dashbox/.style={draw, dashed, rounded corners=5pt, inner sep=8pt}
]
\node[client] (c1) at (0, 0)    {Sumbee\\Server};
\node[client] (c2) at (0,-1.6)  {3rd Party App};
\node[client] (c3) at (0,-3.2)  {Developer\\(curl/Postman)};

\node[dashbox, minimum width=4.2cm, minimum height=5cm] (gwbox) at (4.8,-1.6) {};
\node[font=\footnotesize\bfseries] at (4.8, 0.8) {FastAPI Gateway};
\node[ep] (e1) at (4.8, 0.2)  {POST /analyze/ner};
\node[ep] (e2) at (4.8,-0.8)  {POST /analyze/sentiment};
\node[ep] (e3) at (4.8,-1.8)  {POST /analyze/topics};
\node[ep] (e4) at (4.8,-2.8)  {POST /analyze/full};
\node[ep] (e5) at (4.8,-3.8)  {GET /kb/categories};

\node[dashbox, minimum width=3cm, minimum height=5cm] (nlpbox) at (9.5,-1.6) {};
\node[font=\footnotesize\bfseries] at (9.5, 0.8) {NLP Core};
\node[mod] (m1) at (9.5, 0.2)  {NER Engine};
\node[mod] (m2) at (9.5,-0.8)  {Sentiment};
\node[mod] (m3) at (9.5,-1.8)  {BERTopic};
\node[mod] (m4) at (9.5,-2.8)  {Predefined};
\node[mod] (m5) at (9.5,-3.8)  {KB Manager};

\draw[arrow] (c1)--(e1); \draw[arrow] (c2)--(e3); \draw[arrow] (c3)--(e4);
\draw[arrow] (e1)--(m1); \draw[arrow] (e2)--(m2); \draw[arrow] (e3)--(m3);
\draw[arrow] (e4)--(m1); \draw[arrow] (e4.east) -- ++(0.3,0) |- (m2.west);
\draw[arrow] (e5)--(m5);
\end{tikzpicture}}
\caption{Public REST API-ийн архитектур}
\label{fig:api_arch}
\end{figure}

\begin{table}[h!]
\centering
\renewcommand{\arraystretch}{1.4}
\begin{tabular}{|c|l|p{3.5cm}|p{4.5cm}|}
\hline
\textbf{Method} & \textbf{Endpoint} & \textbf{Input} & \textbf{Output} \\ \hline
POST & \texttt{/analyze/ner} & \texttt{\{texts: [str]\}} & Entity жагсаалт, байрлал, ангилал \\ \hline
POST & \texttt{/analyze/sentiment} & \texttt{\{texts: [str]\}} & Positive/Negative/Neutral оноо \\ \hline
POST & \texttt{/analyze/topics} & \texttt{\{texts, method\}} & Сэдэв, keyword-ууд, хувь \\ \hline
POST & \texttt{/analyze/full} & \texttt{\{texts: [str]\}} & NER + Sentiment + Topic + Network \\ \hline
GET  & \texttt{/kb/categories} & --- & Сэдвийн ангиллын мод бүтэц \\ \hline
\end{tabular}
\caption{Public REST API-ийн үндсэн endpoint-ууд}
\label{tab:api_endpoints}
\end{table}

\subsection{Sumbee.mn системтэй интеграцийн архитектур}

Sumbee.mn интеграцийн хувилбарт NLP pipeline нь Sumbee-ийн Processing Server-т шинэ модуль болгон нэвтэрнэ (зураг~\ref{fig:sumbee_integration}). Одоогийн rule-based боловсруулалтын дараа NLP Pipeline модуль ажиллах ба үр дүнг боловсруулсан өгөгдлийн санд хадгална.

\begin{figure}[h!]
\centering
\resizebox{\linewidth}{!}{%
\begin{tikzpicture}[align=center,
    sbox/.style={rectangle, draw, rounded corners=3pt,
        minimum width=2.5cm, minimum height=0.9cm, text centered, font=\small},
    nlp/.style={rectangle, draw=green!60!black, thick, fill=green!8,
        rounded corners=3pt, minimum width=2.0cm, minimum height=0.75cm,
        text centered, font=\footnotesize},
    arrow/.style={->, >=stealth, thick},
    dashbox/.style={draw=green!60!black, dashed, rounded corners=6pt,
        inner sep=10pt, thick}
]
\node[sbox, fill=blue!10]   (craw) at (0,0)    {Crawler\\Servers};
\node[sbox, fill=blue!10]   (raw)  at (3.2,0)  {Raw\\Storage};
\node[sbox, fill=orange!10] (proc) at (6.5,0)  {Processing\\Server};
\node[sbox, fill=cyan!10]   (pdb)  at (10.5,0) {Processed\\Database};
\node[sbox, fill=gray!10]   (web)  at (13.5,0) {Web App};

\node[dashbox, minimum width=7cm, minimum height=3.2cm] (nlpmodule) at (6.5,-3) {};
\node[font=\small\bfseries, text=green!50!black] at (6.5,-1.3)
    {Шинэ NLP Pipeline модуль};

\node[nlp] (prep)  at (3.8,-2.5) {Mongolian\\Preprocessor};
\node[nlp] (ner)   at (5.8,-2.5) {BERT NER};
\node[nlp] (sent)  at (7.8,-2.5) {Sentiment};
\node[nlp] (topic) at (5.0,-3.7) {BERTopic};
\node[nlp] (net)   at (8.0,-3.7) {Network\\Analyzer};

\draw[arrow] (craw)--(raw);
\draw[arrow] (raw)--(proc);
\draw[arrow] (proc)--(pdb);
\draw[arrow] (pdb)--(web);

\draw[arrow] (proc.south) -- ++(0,-0.4) -| (prep.north);
\draw[arrow] (prep)--(ner);
\draw[arrow] (ner)--(sent);
\draw[arrow] (ner.south)  -- (topic.north);
\draw[arrow] (sent.south) -- (net.north);
\draw[arrow] (topic.east) -- (net.west);
\draw[arrow] (net.north east) -- ++(0.5,0.5) |- (pdb.south);

\node[font=\footnotesize,text=green!50!black] at (6.5,-4.6)
    {NLP шинжилгээний үр дүн → Processed DB};
\end{tikzpicture}}
\caption{Sumbee.mn системтэй интеграцийн архитектур}
\label{fig:sumbee_integration}
\end{figure}

% ================================================================
\section{Системийн ерөнхий Hexagonal архитектур}
% ================================================================

\textbf{Hexagonal Architecture} (Ports \& Adapters) загварын үндсэн санаа нь: \textit{``Application логик нь гадны технологи болон framework-ээс бүрэн тусгаарлагдсан байх''}. Манай системд гурван давхарга байна:

\begin{itemize}
    \item \textbf{Domain Core:} NER, Sentiment, Topic Modeling, Network Analysis-ийн бизнес логик. FastAPI, SQLite, Next.js-ийн талаар мэддэггүй.
    \item \textbf{Ports:} Domain Core-д орох/гарах хийсвэр интерфейс (\texttt{ITextAnalyzer}, \texttt{ITopicModeler}, \texttt{IResultStore}).
    \item \textbf{Adapters:} Портуудыг бодит технологиор хэрэгжүүлнэ: FastAPI Router Adapter, SQLite Adapter, HuggingFace Model Adapter.
\end{itemize}

\begin{figure}[h!]
\centering
\resizebox{0.80\linewidth}{!}{%
\begin{tikzpicture}[align=center,
    core/.style={rectangle, draw=green!60!black, very thick, fill=green!8,
        rounded corners=6pt, minimum width=5.5cm, minimum height=4.0cm},
    port/.style={rectangle, draw=blue!60, thick, fill=blue!8,
        rounded corners=3pt, minimum width=1.8cm, minimum height=0.7cm,
        text centered, font=\footnotesize},
    adap/.style={rectangle, draw=orange!70, thick, fill=orange!10,
        rounded corners=3pt, minimum width=2.2cm, minimum height=0.75cm,
        text centered, font=\footnotesize},
    darrow/.style={<->, >=stealth, thick, draw=gray!60}
]
\node[core] (core) at (0,0) {};
\node[font=\small\bfseries, text=green!50!black] at (0, 1.4) {NLP Domain Core};
\node[font=\footnotesize] at (0, 0.7) {NER Engine};
\node[font=\footnotesize] at (0, 0.2) {Sentiment Analyzer};
\node[font=\footnotesize] at (0,-0.3) {Topic Modeler};
\node[font=\footnotesize] at (0,-0.8) {Network Analyzer};
\node[font=\footnotesize] at (0,-1.3) {KB Manager};

\node[port] (ip1) at (-4.5, 1.0) {Input Port\\ITextAnalyzer};
\node[port] (ip2) at (-4.5,-0.3) {Input Port\\ITopicModeler};
\node[port] (ip3) at (-4.5,-1.6) {Input Port\\IKBManager};
\node[port] (op1) at ( 4.5, 0.7) {Output Port\\IResultStore};
\node[port] (op2) at ( 4.5,-0.8) {Output Port\\IModelLoader};

\node[adap] (a1) at (-7.8, 1.0) {FastAPI\\Router};
\node[adap] (a2) at (-7.8,-0.3) {CSV File\\Adapter};
\node[adap] (a3) at (-7.8,-1.6) {Admin REST\\Adapter};
\node[adap] (b1) at ( 7.8, 0.7) {SQLite\\Adapter};
\node[adap] (b2) at ( 7.8,-0.8) {HuggingFace\\Adapter};

\draw[darrow] (a1)--(ip1); \draw[darrow] (a2)--(ip2); \draw[darrow] (a3)--(ip3);
\draw[darrow] (ip1)--(-2.75, 1.0);
\draw[darrow] (ip2)--(-2.75,-0.3);
\draw[darrow] (ip3)--(-2.75,-1.6);
\draw[darrow] (2.75, 0.7)--(op1);
\draw[darrow] (2.75,-0.8)--(op2);
\draw[darrow] (op1)--(b1); \draw[darrow] (op2)--(b2);

\node[font=\footnotesize, text=blue!60]   at (-4.5, 2.0) {Input Ports};
\node[font=\footnotesize, text=blue!60]   at ( 4.5, 2.0) {Output Ports};
\node[font=\footnotesize, text=orange!70] at (-7.8, 2.0) {Adapters (In)};
\node[font=\footnotesize, text=orange!70] at ( 7.8, 2.0) {Adapters (Out)};
\end{tikzpicture}}
\caption{Hexagonal Architecture (Ports \& Adapters) зохиомж}
\label{fig:hexagonal}
\end{figure}

% ================================================================
\section{NLP Pipeline-ийн мэдээллийн урсгалын диаграмм}
% ================================================================

Системийн мэдээллийн урсгал нь зураг~\ref{fig:nlp_pipeline}-д үзүүлсэн дарааллаар явагддаг. Input текст preprocessing-оор дамжин NER болон Sentiment Analysis зэрэгцэн ажиллаж, үр дүн нь Topic Modeling болон Network Analysis-д ашиглагдана.

\begin{figure}[h!]
\centering
\resizebox{0.60\linewidth}{!}{%
\begin{tikzpicture}[align=center,
    io/.style={trapezium, trapezium left angle=70, trapezium right angle=110,
        draw, fill=gray!15, minimum width=3cm, minimum height=0.75cm,
        text centered, font=\footnotesize},
    proc/.style={rectangle, draw, rounded corners=3pt, fill=blue!10,
        minimum width=3.2cm, minimum height=0.85cm, text centered, font=\footnotesize},
    dec/.style={diamond, draw, fill=yellow!20, minimum width=1.8cm,
        minimum height=0.9cm, text centered, font=\footnotesize, aspect=2.5},
    res/.style={rectangle, draw, rounded corners=3pt, fill=green!12,
        minimum width=3.2cm, minimum height=0.75cm, text centered, font=\footnotesize},
    arrow/.style={->, >=stealth, thick},
    node distance=0.55cm
]
\node[io]   (inp)   at (0, 0)    {CSV / Текст оруулах};
\node[proc] (prep)  at (0,-1.6)  {Preprocessing};
\node[dec]  (val)   at (0,-3.2)  {Хүчинтэй?};
\node[proc] (ner)   at (-2.5,-5) {NER Engine\\(BERT)};
\node[proc] (sent)  at ( 2.5,-5) {Sentiment\\(XLM-R)};
\node[proc] (topic) at (0,-6.8)  {Topic Modeling\\(BERTopic + Predefined)};
\node[proc] (net)   at (0,-8.4)  {Network Analysis};
\node[proc] (agg)   at (0,-10.0) {Үр дүн нэгтгэх};
\node[res]  (cache) at (3.5,-10.0) {Cache (SQLite)};
\node[io]   (out)   at (0,-11.5) {Dashboard / API хариу};

\draw[arrow] (inp)  -- (prep);
\draw[arrow] (prep) -- (val);
\draw[arrow] (val)  -- node[right, font=\footnotesize]{Тийм} (0,-4.1);
\draw[arrow] (val.east) -- ++(1.5,0) node[right, font=\footnotesize]{Үгүй: алдаа};
\draw[arrow] (0,-4.1) -| (ner.north);
\draw[arrow] (0,-4.1) -| (sent.north);
\draw[arrow] (ner.south)  |- (topic.west);
\draw[arrow] (sent.south) |- (topic.east);
\draw[arrow] (topic) -- (net);
\draw[arrow] (net)   -- (agg);
\draw[arrow] (agg)   -- (cache);
\draw[arrow] (agg)   -- (out);
\end{tikzpicture}}
\caption{NLP Pipeline-ийн мэдээллийн урсгалын диаграмм}
\label{fig:nlp_pipeline}
\end{figure}

% ================================================================
\section{Дарааллын диаграмм}
% ================================================================

Зураг~\ref{fig:sequence}-д ``POST /analyze/full'' API хүсэлт илгээхэд системийн дотоод дарааллыг үзүүллээ.

\begin{figure}[h!]
\centering
\resizebox{\linewidth}{!}{%
\begin{tikzpicture}[
    lhdr/.style={rectangle, draw, fill=gray!10, minimum width=2.2cm,
        minimum height=0.6cm, text centered, font=\footnotesize},
    msg/.style={->, >=stealth, thick},
    ret/.style={->, >=stealth, thick, dashed},
    lbl/.style={font=\footnotesize, midway, above}
]
\node[lhdr] at (0,0)    {Хэрэглэгч};
\node[lhdr] at (3.5,0)  {Next.js};
\node[lhdr] at (7,0)    {FastAPI};
\node[lhdr] at (10.5,0) {NLP Core};
\node[lhdr] at (14,0)   {SQLite};

\foreach \x in {0,3.5,7,10.5,14}
    \draw[dashed, gray] (\x,-0.3) -- (\x,-10.5);

\draw[msg] (0,-1.0)  -- (3.5,-1.0)  node[lbl]{1. CSV оруулах};
\draw[msg] (3.5,-1.8)-- (7,-1.8)    node[lbl]{2. POST /analyze/full};
\draw[msg] (7,-2.6)  -- (10.5,-2.6) node[lbl]{3. preprocess(texts)};
\draw[ret] (10.5,-3.2)-- (7,-3.2)   node[lbl]{4. cleaned\_texts};
\draw[msg] (7,-4.0)  -- (10.5,-4.0) node[lbl]{5. ner\_extract()};
\draw[msg] (7,-4.8)  -- (10.5,-4.8) node[lbl]{6. sentiment\_predict()};
\draw[msg] (7,-5.6)  -- (10.5,-5.6) node[lbl]{7. topic\_model()};
\draw[msg] (7,-6.4)  -- (10.5,-6.4) node[lbl]{8. network\_analyze()};
\draw[msg] (10.5,-7.2)-- (14,-7.2)  node[lbl]{9. save\_cache()};
\draw[ret] (14,-7.9) -- (10.5,-7.9) node[lbl]{10. cache\_id};
\draw[ret] (10.5,-8.7)-- (7,-8.7)   node[lbl]{11. result JSON};
\draw[ret] (7,-9.5)  -- (3.5,-9.5)  node[lbl]{12. HTTP 200 JSON};
\draw[ret] (3.5,-10.2)-- (0,-10.2)  node[lbl]{13. Dashboard дүрслэх};
\end{tikzpicture}}
\caption{``POST /analyze/full'' хүсэлтийн дарааллын диаграмм}
\label{fig:sequence}
\end{figure}

% ================================================================
\section{Өгөгдлийн сангийн зохиомж}
% ================================================================

Системийн мэдлэгийн санг \textbf{SQLite} өгөгдлийн санд (\texttt{knowledge.db}) хадгална. SQLite-ийг сонгосон шалтгаан: суулгалт шаардлагагүй, нэг файл, хялбар backup, 10K+ бичлэгт хангалттай гүйцэтгэлтэй.

\begin{figure}[h!]
\centering
\resizebox{0.82\linewidth}{!}{%
\begin{tikzpicture}[
    tbl/.style={rectangle, draw, fill=blue!8, rounded corners=4pt,
        minimum width=2.8cm, minimum height=0.9cm,
        text centered, font=\small\bfseries},
    rel/.style={->, >=stealth, thick, draw=gray!70},
    node distance=1.2cm
]
\node[tbl] (sw)    at (0, 0)    {stopwords};
\node[tbl] (cat)   at (5.5, 0)  {categories};
\node[tbl] (lbl)   at (5.5,-3)  {labels};
\node[tbl] (hist)  at (0,-3)    {analysis\_history};
\node[tbl] (seeds) at (11, 0)   {topic\_seeds};

\draw[rel] (cat.south) -- node[right,font=\footnotesize]{1:N} (lbl.north);
\draw[rel] (cat.east)  -- node[above,font=\footnotesize]{1:N} (seeds.west);
\draw[rel] (cat.north west) to[bend left=50]
    node[above,font=\footnotesize]{self-ref} (cat.north east);

\node[font=\footnotesize, text=gray, align=left] at (0, -0.7)
    {id, word, language};
\node[font=\footnotesize, text=gray, align=left] at (5.5, -0.7)
    {id, name, parent\_id, embedding};
\node[font=\footnotesize, text=gray, align=left] at (5.5, -3.7)
    {id, name, color, category\_id};
\node[font=\footnotesize, text=gray, align=left] at (0, -3.7)
    {id, input\_hash, result\_json};
\node[font=\footnotesize, text=gray, align=left] at (11, -0.7)
    {id, topic\_name, category\_id};
\end{tikzpicture}}
\caption{Knowledge Base өгөгдлийн сангийн ER диаграмм}
\label{fig:er_diagram}
\end{figure}

\begin{table}[h!]
\centering
\renewcommand{\arraystretch}{1.4}
\begin{tabular}{|l|p{3.5cm}|p{5.5cm}|}
\hline
\textbf{Хүснэгт} & \textbf{Гол талбарууд} & \textbf{Зориулалт} \\ \hline
\texttt{stopwords} & id, word, language & Монгол/Англи зогсох үгс. Preprocessing-д ашиглагдана. \\ \hline
\texttt{categories} & id, name, parent\_id, embedding & Сэдвийн ангиллын мод бүтэц. Predefined topic-ийн embedding. \\ \hline
\texttt{labels} & id, name, color, category\_id & NER болон Topic шошгоны дүрслэлийн тохиргоо. \\ \hline
\texttt{analysis\_history} & id, input\_hash, result\_json & Анализын кэш. Ижил input дахин ирвэл кэшийг буцаана. \\ \hline
\texttt{topic\_seeds} & id, topic\_name, embedding & Predefined topic-ийн жишиг өгүүлбэр болон embedding. \\ \hline
\end{tabular}
\caption{Knowledge Base хүснэгтүүдийн бүтэц ба зориулалт}
\label{tab:db_tables}
\end{table}

% ================================================================
\section{NLP Загваруудын сонголт ба харьцуулалт}
% ================================================================

\subsection{NER загварын харьцуулалт}

NER даалгаварт зориулж гурван нэр дэвшигч загварыг харьцуулан судлав.

\begin{table}[h!]
\centering
\renewcommand{\arraystretch}{1.5}
\begin{tabular}{|p{3.8cm}|c|c|c|c|c|c|}
\hline
\textbf{Загварын нэр} & \textbf{Хэл} & \textbf{Ангилал} & \textbf{Монгол} & \textbf{GPU (ms)} & \textbf{CPU (ms)} & \textbf{Сонгосон} \\ \hline
\texttt{Davlan/bert-multilingual\--ner-hrl} & 10 & PER, ORG, LOC & Тийм & 20--60 & 300--600 & \checkmark \\ \hline
\texttt{spaCy xx\_ent\_wiki\_sm} & Олон & PER, ORG, LOC & Үгүй & 5 & 10 & $\times$ \\ \hline
\texttt{flair/ner-multi} & Олон & PER,ORG,LOC,MISC & Хязгаарлагдмал & 100 & 500 & $\times$ \\ \hline
\end{tabular}
\caption{NER загваруудын харьцуулалт}
\label{tab:ner_comparison}
\end{table}

\texttt{Davlan/bert-base-multilingual-cased-ner-hrl} загварыг сонгосон үндэслэл: агглютинатив бүтэцтэй хэлүүдийн (Turkish, Finnish гэх мэт) өгөгдөл дээр pretrain хийгдсэн multilingual embedding нь Монгол нэр таних чадвартай. spaCy xx загвар Монгол хэлийг дэмжихгүй тул хасагдсан.

\subsection{Topic Modeling загварын харьцуулалт}

\begin{table}[h!]
\centering
\renewcommand{\arraystretch}{1.5}
\begin{tabular}{|p{3.5cm}|c|c|c|}
\hline
\textbf{Хэмжүүр} & \textbf{LDA} & \textbf{BERTopic} & \textbf{Predefined Class.} \\ \hline
Сургалтын датасет & Шаардлагатай & Шаардлагагүй & Шаардлагагүй \\ \hline
Topic тоо тогтоох & Гараар ($K$) & Автомат & Тогтсон \\ \hline
Монгол текст & Дунд & Сайн & Сайн \\ \hline
Inference хурд & Хурдан & Дунд (200ms) & Маш хурдан (5ms) \\ \hline
Шинэ сэдэв илрүүлэх & Тийм & Тийм & Үгүй \\ \hline
Тайлбарлагдах байдал & Өндөр & Дунд & Өндөр \\ \hline
Сонгосон & $\times$ & \checkmark & \checkmark \\ \hline
\end{tabular}
\caption{Topic Modeling аргуудын харьцуулалт}
\label{tab:topic_comparison}
\end{table}

BERTopic болон Predefined Classification-ийг хоёуланг нь сонгосон шалтгаан нь: BERTopic нь урьдчилан мэдэгдэхгүй шинэ сэдэв илрүүлэхэд давуу, харин Predefined Classification нь тодорхой бизнесийн ангиллыг маш хурдан (1--5 ms) таньдаг. Хоёр аргыг хослуулан ашиглах нь хамгийн оновчтой үр дүн өгдөг.

% ================================================================
\section{Дүгнэлт}
% ================================================================

Системийн шинжилгээ ба зохиомжийн бүлгийн дүгнэлтэнд: нийт 10 функциональ болон 8 функциональ бус шаардлагыг тодорхойлж, Use Case диаграммаар гурван актор, 10 хэрэглээний тохиолдлыг дүрсэллээ. Hexagonal Architecture нь NLP Domain Core-г технологиос тусгаарлаж, Demo Web, Public API, Sumbee интеграци гэсэн гурван хувилбарыг нэг кодоос хэрэгжүүлэх боломжийг бүрдүүлж байна. Өгөгдлийн сангийн 5 хүснэгтийн зохиомж болон NLP загваруудын харьцуулалтаас \texttt{Davlan/bert-multilingual-ner-hrl} болон BERTopic+Predefined хосолсон арга нь хамгийн тохиромжтой шийдэл болохыг тогтоов. Дараагийн бүлэгт эдгээр зохиомжийн хэрэгжүүлэлт болон туршилтын үр дүнг танилцуулна.

\chapter{Хэрэгжүүлэлт ба туршилт}


% ================================================================
\section{Ашигласан технологиуд}
% ================================================================

\subsection{Backend технологиуд}

\begin{table}[h!]
\centering
\renewcommand{\arraystretch}{1.4}
\begin{tabular}{|l|l|p{6cm}|}
\hline
\textbf{Технологи} & \textbf{Хувилбар} & \textbf{Зориулалт} \\ \hline
Python & 3.11 & Үндсэн хэл — NLP core болон API backend \\ \hline
FastAPI & 0.110 & Async REST API framework, OpenAPI автомат \\ \hline
HuggingFace Transformers & 4.40 & BERT NER болон XLM-R Sentiment pipeline \\ \hline
BERTopic & 0.16 & Sentence embedding + HDBSCAN/KMeans кластер \\ \hline
sentence-transformers & 2.7 & Multilingual sentence embedding (paraphrase-mpnet) \\ \hline
NetworkX & 3.3 & Entity co-occurrence граф тооцоолол \\ \hline
SQLite / sqlite3 & 3.45 & Knowledge Base өгөгдлийн сан \\ \hline
Pydantic & 2.7 & API schema validation, dataclass \\ \hline
Docker & 25 & Container-д суурилсан байршуулалт \\ \hline
\end{tabular}
\caption{Backend технологиуд}
\label{tab:backend_tech}
\end{table}

\subsection{Frontend технологиуд}

\begin{table}[h!]
\centering
\renewcommand{\arraystretch}{1.4}
\begin{tabular}{|l|l|p{6cm}|}
\hline
\textbf{Технологи} & \textbf{Хувилбар} & \textbf{Зориулалт} \\ \hline
Next.js & 14 (App Router) & React-д суурилсан SSR/CSR framework \\ \hline
Tailwind CSS & 3.4 & Utility-first CSS, glassmorphic dark-mode \\ \hline
Recharts & 2.12 & Pie chart, Bar chart, Timeline дүрслэл \\ \hline
React Force Graph & 1.44 & D3.js-д суурилсан network graph дүрслэл \\ \hline
TypeScript & 5.4 & Статик тип шалгалт \\ \hline
\end{tabular}
\caption{Frontend технологиуд}
\label{tab:frontend_tech}
\end{table}

% ================================================================
\section{Preprocessing модулийн хэрэгжүүлэлт}
% ================================================================

\subsection{Хоёр горимт preprocessing архитектур}

Монгол сошиал медиа текст боловсруулахад нэг preprocessing горим хангалтгүй байгааг туршилтын явцад тодорхойлов. NER болон Sentiment Analysis-д BERT загвар текстийн бүтэц, том-жижиг үсгийг мэдрэмтгий байдлаар ашигладаг тул \textbf{preprocess\_nlp()} горим цэгтэмдэг, том үсгийг хадгалж, нэрийн хэлбэрийг сэргээдэг. BERTopic-д зориулсан \textbf{preprocess\_tm()} горим харин агуулга бүхий токенуудыг л үлдээхийн тулд агрессив цэвэрлэгээ хийдэг.

\begin{lstlisting}[language=Python, caption={Preprocessing хоёр горим — preprocess\_dual()}, basicstyle=\small\ttfamily, breaklines=true, frame=single, numbers=left]
class Preprocessor:
    def preprocess_nlp(self, text: str) -> str:
        """NER + Sentiment-д: цэгтэмдэг, том үсэг хадгалж,
        нэрийн хэлбэрийг сэргээнэ."""
        text = self.clean_basic(text, replace_url=True)
        text = _capitalize_for_ner(text)   # б.амар → Б.Амар
        text = _restore_names(text)        # Б_Амар → Б.Амар
        return text

    def preprocess_tm(self, text: str) -> str:
        """BERTopic-д: агрессив цэвэрлэгээ, зогсох үгс хасна."""
        text = self.clean_basic(text, replace_url=False)
        text = self.clean_deep(text)      # зөвхөн Кирилл/Latin/тоо
        text = text.lower()
        text = self._remove_stopwords(text)
        return text

    def preprocess_dual(self, text: str) -> Tuple[str, str]:
        """Нэг дуудлагаар хоёр горимын үр дүнг буцаана."""
        return self.preprocess_nlp(text), self.preprocess_tm(text)
\end{lstlisting}

\subsection{Монгол нэр хамгаалах механизм}

Монгол нэрийн ``А.Бат'' хэлбэр нь нийгмийн сүлжээнд маш түгээмэл. Preprocessing дуусахад цэг нь арилж ``А Бат'' болдог тул BERT загвар нэрийг таньдаггүй болж байсан. Энэ асуудлыг нэрийг ``А\_Бат'' хэлбэрт хамгаалж, дараа нь сэргээх механизмаар шийдлээ:

\begin{lstlisting}[language=Python, caption={Нэр хамгаалах regex — дээд болон жижиг үсэг хоёуланд}, basicstyle=\small\ttfamily, breaklines=true, frame=single, numbers=left]
# Дээд үсгийн эхний үсэг: А.Бат-Эрдэнэ → А_Бат-Эрдэнэ
MN_NAME_UPPER = re.compile(
    r"\b([А-ЯӨҮЁ])\.\s*"
    r"([А-Яа-яӨөҮүЁё][а-яөүё]+"
    r"(?:-[А-Яа-яӨөҮүЁё][а-яөүё]+)*)"
)
# Жижиг үсгийн эхний үсэг: б.амар → Б_Амар
MN_NAME_LOWER = re.compile(
    r"\b([а-яөүё])\.\s*([а-яөүё]+(?:-[а-яөүё]+)*)"
)
\end{lstlisting}

\textbf{Туршилтын үр дүн:} Нэр хамгаалах механизм нэмэхээс өмнө ``б.амар улс төрч'' текстэд NER ямар ч PER entity буцааж байгаагүй. Нэмсний дараа ``Б.Амар'' PER entity-г 0.94 итгэлцлийн оноотойгоор таньдаг болов.

\subsection{Preprocessing аргуудын давуу ба сул тал}

\begin{table}[h!]
\centering
\renewcommand{\arraystretch}{1.4}
\begin{tabular}{|p{3.5cm}|p{5cm}|p{5cm}|}
\hline
\textbf{Арга} & \textbf{Давуу тал} & \textbf{Сул тал} \\ \hline
Нэг горимт preprocessing & Энгийн, хялбар & NER болон TM аль нэгд тохиромжгүй \\ \hline
Давхар горимт (dual mode) & Тус бүрийн задач-д оптимал & Нэг текстийг хоёр удаа боловсруулна (2× хугацаа) \\ \hline
Нэр хамгаалах (protect) & Монгол нэрийн таних чадвар +30\% & Regex-ийн нарийвчлал шаардана \\ \hline
\end{tabular}
\caption{Preprocessing аргуудын харьцуулалт}
\label{tab:preprocessing_comparison}
\end{table}

% ================================================================
\section{NER модулийн хэрэгжүүлэлт}
% ================================================================

\subsection{Lazy-load ба batch inference архитектур}

HuggingFace-ийн BERT загвар (~700MB) эхний дуудлагад л ачаалагддаг ``lazy-load'' аргыг хэрэглэв. Batch inference нь олон текстийг нэг дор GPU/CPU-д дамжуулж боловсруулдаг тул нэг-нэгээр дамжуулахаас ~4 дахин хурдан байна.

\begin{lstlisting}[language=Python, caption={NEREngine — lazy-load ба batch inference}, basicstyle=\small\ttfamily, breaklines=true, frame=single, numbers=left]
class NEREngine:
    def __init__(self,
      model_name="Davlan/bert-base-multilingual-cased-ner-hrl"):
        self.model_name = model_name
        self._pipeline = None  # lazy-load: эхний дуудлагад ачаална

    def _load_pipeline(self):
        if self._pipeline is None:
            from transformers import pipeline
            self._pipeline = pipeline(
                "ner",
                model=self.model_name,
                aggregation_strategy="simple",  # B-/I- токен нэгтгэнэ
            )
        return self._pipeline

    def recognize_batch(self, texts: List[str],
                        batch_size: int = 16) -> List[List[EntityResult]]:
        """Batch processing: 16 текстийг нэг GPU forward pass-д."""
        pipe = self._load_pipeline()
        raw_results = pipe(valid_texts, batch_size=batch_size)
        # ... entity_group, score, start, end буцаана
\end{lstlisting}

\subsection{NER туршилтын үр дүн (fine-tune хийхээс өмнө)}

Суурь загвар (\texttt{Davlan/bert-base-multilingual-cased-ner-hrl})-ыг Sumbee.mn-ийн 500 нийтлэл дээр entity-level exact match аргаар үнэлэхэд дараах үр дүн гарав:

\begin{table}[h!]
\centering
\renewcommand{\arraystretch}{1.4}
\begin{tabular}{|l|c|c|c|c|}
\hline
\textbf{Үзүүлэлт} & \textbf{Precision} & \textbf{Recall} & \textbf{F1} & \textbf{Тайлбар} \\ \hline
Суурь загвар (500 нийтлэл) & 0.321 & 0.316 & \textbf{0.318} & Fine-tune хийгдээгүй \\ \hline
\end{tabular}
\caption{NER суурь загварын гүйцэтгэлийн үнэлгээ (eval/evaluate\_tokens.py, 500 нийтлэл)}
\label{tab:ner_results}
\end{table}

\begin{itemize}
    \item \textbf{True Positives:} 226 — зөв олсон entity
    \item \textbf{False Positives:} 478 — буруу таньсан entity
    \item \textbf{False Negatives:} 489 — алдсан entity
\end{itemize}

F1=0.318 нь бодит хэрэглээнд тохирохгүй бага үзүүлэлт юм. Гол шалтгаанууд нь: (1) загвар Монгол хэл дээр тусгайлан fine-tune хийгдээгүй, (2) Монгол нэрийн ``А.Бат-Эрдэнэ'' хэлбэрийг хуваан таньдаг, (3) ``МАН'', ``УИХ'' зэрэг Монгол товчлолыг мэддэггүй.

\textbf{Давуу тал:} Монгол хэлний агглютинатив бүтцийг BERT subword tokenization зохицуулж чаддаг; batch inference нь GPU дээр 20--60 ms латентитэй. \textbf{Сул тал:} Fine-tune хийхгүйгээр Монгол текст дээр F1=0.318 — домайн-специфик нэршил, товчлол алдагддаг.

% ================================================================
\section{NER загварын нарийн тааруулга (Fine-tuning)}
% ================================================================

Суурь загвар (\texttt{Davlan/bert-base-multilingual-cased-ner-hrl}) нь Монгол хэл дээр тусгайлан сургаагүй тул хэд хэдэн чухал дутагдалтай байв. Тухайлбал ``Л.Оюун-Эрдэнэ'' нэрийг ``Л'' болон ``Оюун-Эрдэнэ'' гэж хоёр тусдаа entity болгон хуваадаг, ``МАН'' зэрэг Монгол байгууллагын товчлолыг таньдаггүй, мөн газрын нэрийн нарийн ялгааг сайн ялгадаггүй байсан. Эдгээр асуудлыг шийдэхийн тулд хоёр үе шаттай fine-tuning хийв.

\subsection{Асуудлын тодорхойлолт — Tokenization ба тэмдэглэлийн зөрчил}

BERT загварын WordPiece tokenizer нь текстийг sub-token (дэд-токен) болгон хуваадаг. Монгол хэлний нэрийн ``Д.Батулга'' хэлбэрт энэ нь онцгой асуудал үүсгэдэг:

\begin{lstlisting}[caption={WordPiece tokenization — цэг дээр хуваагдах жишээ}, basicstyle=\small\ttfamily, breaklines=true, frame=single]
Оригинал:    "Д.Батулга"
Токенууд:    ["Д", ".", "Батулга"]     (3 токен)

V1 тэмдэглэл (word-level):
  Д        → B-PER
  .        → -100    ← загвар "." тэмдэгтийг нэрийн дотор
  Батулга  → -100       гэдгийг хэзээ ч суралцдаггүй!

V2 тэмдэглэл (character-offset):
  Д        → B-PER
  .        → I-PER   ← загвар "." нь нэрийн дотор
  Батулга  → I-PER      байж болохыг суралцдаг!
\end{lstlisting}

V1 сургалтын өгөгдөлд зөвхөн эхний дэд-токен тэмдэглэл авч, үлдсэн нь $-100$ (``үл тооцох'') утга авдаг байсан. Үүний улмаас загвар ``\texttt{.}'' тэмдэгтийг нэрийн дотор байж болохыг хэзээ ч суралцаагүй бөгөөд inference хийхэд ``\texttt{.}'' дээр entity-г таслан хуваадаг байв. Ижил асуудал ``\texttt{-}'' (зураас) тэмдэгт дээр ч тохиолддог байсан (жишээ нь ``Оюун-Эрдэнэ'').

\subsection{V1 сургалт — Суурь fine-tuning}

Эхний fine-tuning нь CoNLL форматын (үг-level тэмдэглэл) өгөгдөл ашигласан бөгөөд \texttt{Davlan/bert-base-multilingual-cased-ner-hrl} суурь загвараас эхлэн сургасан. Google Colab-ийн T4 GPU дээр 4 epoch сургасан:

\begin{table}[h!]
\centering
\renewcommand{\arraystretch}{1.4}
\begin{tabular}{|c|c|c|c|c|c|}
\hline
\textbf{Epoch} & \textbf{Training Loss} & \textbf{Val Loss} & \textbf{F1} & \textbf{Precision} & \textbf{Recall} \\ \hline
1 & 0.0875 & 0.1042 & 0.773 & 0.752 & 0.794 \\ \hline
2 & 0.0587 & 0.1027 & 0.801 & 0.773 & 0.830 \\ \hline
3 & 0.0390 & 0.1007 & 0.811 & 0.771 & 0.856 \\ \hline
4 & 0.0259 & 0.1108 & 0.809 & 0.786 & 0.834 \\ \hline
\end{tabular}
\caption{V1 fine-tuning сургалтын үр дүн (CoNLL word-level тэмдэглэл)}
\label{tab:v1_training}
\end{table}

V1 загвар F1=0.811 хүрсэн нь суурь загвараас мэдэгдэхүйц сайжирсан боловч Монгол нэрийн ``А.Бат-Эрдэнэ'' хэлбэрийг зөв таньдаггүй сул талтай байсан. Eval скрипт (\texttt{eval/evaluate\_tokens.py})-ийг Sumbee.mn-ийн 500 нийтлэл дээр ажиллуулахад entity-level F1=0.32 гарсан нь бодит хэрэглээнд тохирохгүй байгааг баталсан.

\subsection{V2 сургалтын өгөгдлийн бэлтгэл}

V2 сургалтын өгөгдлийг бэлтгэхэд дараах алхмуудыг хэрэгжүүлэв:

\subsubsection{Тэмдэглэлийн формат — Character-offset}

Word-level тэмдэглэлийн оронд character-offset (тэмдэгт-хэсгийн) тэмдэглэлийг ашиглав. Энэ форматад entity-ийн эхлэл ба төгсгөлийн тэмдэгтийн байрлалыг (start, end) тодорхойлж, тухайн хэсэгт багтах \textbf{бүх} дэд-токенд BIO тэмдэглэл өгдөг:

\begin{lstlisting}[caption={Character-offset JSON формат}, basicstyle=\small\ttfamily, breaklines=true, frame=single]
{
  "text": "Ерөнхий сайд Л.Оюун-Эрдэнэ мэдэгдэл хийлээ",
  "labels": [[13, 27, "PER"]]
}
// Тэмдэгт байрлал [13:27] = "Л.Оюун-Эрдэнэ"
// Бүх дэд-токен (Л, ., Оюун, -, Эрдэнэ) B-PER/I-PER авна
\end{lstlisting}

\subsubsection{Өгөгдлийн эх сурвалжууд}

Сургалтын өгөгдлийг олон эх сурвалжаас нэгтгэн үүсгэв:

\begin{table}[h!]
\centering
\renewcommand{\arraystretch}{1.4}
\begin{tabular}{|l|c|p{5.5cm}|}
\hline
\textbf{Эх сурвалж} & \textbf{Тоо} & \textbf{Тайлбар} \\ \hline
NER\_v1.0.json.gz & 10,162 & Монгол NER корпус — гар тэмдэглэлт, PER/ORG/LOC/MISC \\ \hline
Монгол нэрс & 120+ & Түгээмэл эрэгтэй, эмэгтэй нэрс \\ \hline
Улс төрчид & 30+ & Одоогийн бодит нэр, үсгийн товчлолтой (Л.Оюун-Эрдэнэ, У.Хүрэлсүх г.м.) \\ \hline
Товчлолууд & 526 & Монгол байгууллагын товчлол ба бүтэн нэр (УИХ, МАН г.м.) \\ \hline
Засгийн газрын байгууллага & 30+ & Яам, агентлагийн бүтэн нэр \\ \hline
Компани / банк / хэвлэл мэдээлэл & 50+ & Голомт банк, Мобиком, Өрнөдөс ТВ г.м. \\ \hline
Газар нэр (level 2--5) & 530+ & 22 аймаг, 26 сум, 89 хороо, 142 хороолол \\ \hline
\end{tabular}
\caption{V2 сургалтын өгөгдлийн эх сурвалжууд}
\label{tab:v2_data_sources}
\end{table}

\subsubsection{Өгөгдлийн үүсгэлтийн онцлог шийдлүүд}

\begin{enumerate}
    \item \textbf{Цэгтэй нэрийн загвар:} ``Д.Батулга'', ``Л.Оюун-Эрдэнэ'' зэрэг Монгол нэрийн стандарт хэлбэрийг олон тооны template өгүүлбэрт оруулж, загварыг цэгийн дотор entity үргэлжилдгийг суралцуулав.
    \item \textbf{Залгаварт нэрийн хэлбэр:} Монгол хэлний ``-ын'', ``-ийн'', ``-д'', ``-г'' гэсэн нөхцөл залгасан нэрүүдийг (``Батулгыг'', ``Хүрэлсүхийн'') тусгайлан оруулав.
    \item \textbf{Товчлол + бүтэн нэр хослол:} ``МАН'' товчлол болон ``Монгол Ардын Нам'' бүтэн нэрийг аль алиныг нь ORG entity-д сургав.
    \item \textbf{Хольсон entity өгүүлбэрүүд:} Нэг өгүүлбэрт PER + ORG + LOC гурвуулаа орсон жишээнүүд оруулж, загварыг олон entity-тэй контекстэд ажиллахад сургав.
\end{enumerate}

\subsubsection{Нэгтгэсэн өгөгдлийн статистик}

\begin{table}[h!]
\centering
\renewcommand{\arraystretch}{1.4}
\begin{tabular}{|l|c|}
\hline
\textbf{Үзүүлэлт} & \textbf{Утга} \\ \hline
Нийт өгүүлбэр & 20,696 \\ \hline
LOC тэмдэглэл & 11,246 \\ \hline
ORG тэмдэглэл & 10,933 \\ \hline
PER тэмдэглэл & 9,603 \\ \hline
MISC тэмдэглэл (O болгон хөрвүүлсэн) & 2,839 \\ \hline
Цэгтэй нэр entity & 6,862 \\ \hline
Зураастай нэр entity & 1,185 \\ \hline
Нөхцөл залгасан нэр entity & 8,308 \\ \hline
Валидацийн алдаа & 0 \\ \hline
\end{tabular}
\caption{V2 сургалтын өгөгдлийн статистик (train\_v2\_merged.jsonl)}
\label{tab:v2_data_stats}
\end{table}

\subsection{V2 сургалтын тохиргоо ба явц}

V2 fine-tuning нь V1 загвараас үргэлжлүүлэн сургасан (continued fine-tuning). Суурь загвараас шинээр эхлэхийн оронд V1-д суралцсан мэдлэг дээр нэмж сургав.

\subsubsection{Сургалтын тохиргоо}

\begin{table}[h!]
\centering
\renewcommand{\arraystretch}{1.4}
\small
\begin{tabular}{|p{3cm}|p{2.8cm}|p{3.5cm}|p{4.5cm}|}
\hline
\textbf{Параметр} & \textbf{V1} & \textbf{V2} & \textbf{Тайлбар} \\ \hline
Суурь загвар & Davlan BERT & Nomio4640/ ner-mongolian (V1) & V2 нь V1-ээс үргэлжлүүлсэн \\ \hline
Learning rate & 3e-5 & 2e-5 & Үргэлжлүүлэн сургахад бага LR \\ \hline
Epoch & 4 & 5 & Шинэ өгөгдлийн хэмжээнд тохируулсан \\ \hline
Batch size & 16 & 16 & T4 GPU-ийн санах ойд тохирсон \\ \hline
Warmup steps & 200 & 300 & Илүү их өгөгдөлтэй тул нэмсэн \\ \hline
Weight decay & 0.01 & 0.01 & Ижил регуляризаци \\ \hline
FP16 & Тийм & Тийм & Хагас нарийвчлалтай сургалт \\ \hline
Тэмдэглэлийн формат & CoNLL word-level & JSON character-offset & Бүх дэд-токен тэмдэглэл авна \\ \hline
Сургалтын өгүүлбэр & $\sim$8,000 & 20,696 & 2.5 дахин их өгөгдөл \\ \hline
Валидацийн өгүүлбэр & 915 & 915 & Ижил eval set \\ \hline
\end{tabular}
\normalsize
\caption{V1 ба V2 сургалтын тохиргооны харьцуулалт}
\label{tab:v1_v2_config}
\end{table}

\subsubsection{Character-offset tokenization алгоритм}

V2 сургалтын гол шинэлэг шийдэл нь character-offset-д суурилсан токенчлолын тааруулга юм. Уг алгоритм нь HuggingFace tokenizer-ийн \texttt{return\_offsets\_mapping} функцийг ашиглан токен бүрийн анхны текст дэх тэмдэгтийн байрлалыг авч, entity span-д харьцуулдаг:

\begin{lstlisting}[language=Python, caption={Character-offset tokenization тааруулга}, basicstyle=\small\ttfamily, breaklines=true, frame=single, numbers=left]
def tokenize_json_entry(entry):
    text = entry['text']
    label_spans = entry['labels']  # [[start, end, label], ...]
    encoding = tokenizer(text, truncation=True, max_length=512,
                         return_offsets_mapping=True)
    offset_mapping = encoding['offset_mapping']
    token_labels = []
    for idx, (tok_start, tok_end) in enumerate(offset_mapping):
        if tok_start == 0 and tok_end == 0:
            token_labels.append(-100)  # [CLS], [SEP] тусгай токен
            continue
        label = "O"
        for span_start, span_end, span_label in label_spans:
            if tok_start >= span_start and tok_end <= span_end:
                # Өмнөх токен мөн entity дотор бол I-, үгүй бол B-
                is_first = True
                if idx > 0:
                    prev_s, prev_e = offset_mapping[idx - 1]
                    if prev_s >= span_start and prev_e <= span_end:
                        is_first = False
                label = f"B-{span_label}" if is_first else f"I-{span_label}"
                break
        token_labels.append(label2id.get(label, 0))
    encoding['labels'] = token_labels
    return encoding
\end{lstlisting}

Энэ алгоритмын гол давуу тал нь entity span дотор орших \textbf{бүх} тэмдэгт — ``\texttt{.}'', ``\texttt{-}'', дэд-үг (subword) зэрэг — B-/I- тэмдэглэл авдаг. Word-level аргаас ялгаатай нь $-100$ (үл тооцох) утга зөвхөн тусгай токен ([CLS], [SEP], [PAD]) болон entity-ийн гаднах padding-д хэрэглэгддэг.

\subsection{V2 сургалтын үр дүн}

Google Colab T4 GPU дээр 5 epoch, нийт 46 минутын сургалт хийв:

\begin{table}[h!]
\centering
\renewcommand{\arraystretch}{1.4}
\begin{tabular}{|c|c|c|c|c|c|}
\hline
\textbf{Epoch} & \textbf{Training Loss} & \textbf{Val Loss} & \textbf{F1} & \textbf{Precision} & \textbf{Recall} \\ \hline
1 & 0.0388 & 0.2818 & 0.797 & 0.732 & 0.876 \\ \hline
2 & 0.0195 & 0.3036 & 0.822 & 0.756 & 0.901 \\ \hline
3 & 0.0147 & 0.3283 & 0.833 & 0.763 & 0.917 \\ \hline
4 & 0.0083 & 0.3422 & 0.843 & 0.776 & 0.924 \\ \hline
\textbf{5} & \textbf{0.0048} & \textbf{0.3629} & \textbf{0.845} & \textbf{0.778} & \textbf{0.925} \\ \hline
\end{tabular}
\caption{V2 fine-tuning сургалтын үр дүн (character-offset тэмдэглэл)}
\label{tab:v2_training}
\end{table}

\subsection{Сургалтын явцын шинжилгээ}

\subsubsection{Training Loss — Суралцалтын явц}

Training loss нь epoch бүрт тасралтгүй буурч (0.0388 → 0.0048) загвар сургалтын өгөгдлийг маш сайн суралцаж байгааг харуулав. Epoch 5-д loss 0.0048 хүрсэн нь загвар сургалтын өгүүлбэрүүдийн entity-г бараг бүрэн зөв таньж байгааг илтгэнэ.

\subsubsection{Validation Loss — Ерөнхийлөх чадвар}

Validation loss нь epoch 1-ийн 0.2818-аас epoch 5-д 0.3629 болж нэмэгдсэн. Энэ нь дараах шалтгаантай:

\begin{enumerate}
    \item \textbf{Тархалтын зөрүү (Distribution mismatch):} Сургалтын өгөгдөл нь үүсгэсэн (generated) өгүүлбэрүүд, харин валидацийн өгөгдөл нь бодит текстээс авсан CoNLL формат. Хоёр өгөгдлийн хэв маяг ялгаатай.
    \item \textbf{Тэмдэглэлийн формат зөрүү:} Сургалтын өгөгдөлд ``\texttt{.}'' нь I-PER тэмдэглэл авдаг, харин валидацийн өгөгдөлд $-100$ авдаг. Загвар шинэ зүйл суралцаж байгаа ч хуучин eval set үүнийг бүрэн хэмжээгүй.
    \item \textbf{Хэмжээний зөрүү:} Сургалтын 20,696 ба валидацийн 915 өгүүлбэрийн харьцаа (96:4) нь стандарт 80:20 харьцаанаас зөрүүтэй.
\end{enumerate}

Гэвч F1 оноо тасралтгүй өссөн нь validation loss-оос илүү чухал үзүүлэлт юм — загвар зөв хариулт гаргаж байгаа ч итгэлцлийн тархалт (confidence distribution) хуучин eval set-тэй бүрэн нийцэхгүй байна гэсэн үг.

\subsubsection{Precision ба Recall-ийн тэнцвэр}

\begin{itemize}
    \item \textbf{Recall} (0.925) нь \textbf{Precision} (0.778)-ээс мэдэгдэхүйц өндөр. Энэ нь загвар бодит entity-ийн 92.5\%-ийг олдог боловч олсон entity-ийн 22.2\% нь буруу (false positive) байна гэсэн үг.
    \item Сошиал мониторингийн хэрэглээнд recall өндөр байх нь чухал — entity орхигдуулахаас илүү нэмэлт entity-г дараа нь шүүх нь хялбар.
\end{itemize}

\subsection{V1 ба V2 загварын харьцуулалт}

\begin{table}[h!]
\centering
\renewcommand{\arraystretch}{1.4}
\begin{tabular}{|l|c|c|c|}
\hline
\textbf{Үзүүлэлт} & \textbf{Суурь загвар} & \textbf{V1} & \textbf{V2} \\ \hline
F1 Score & 0.318 & 0.811 & \textbf{0.845} \\ \hline
Precision & 0.321 & 0.786 & 0.778 \\ \hline
Recall & 0.316 & 0.834 & \textbf{0.925} \\ \hline
``Д.Батулга'' таних & 2 entity & 2 entity & \textbf{1 entity} \\ \hline
``Л.Оюун-Эрдэнэ'' таних & 3 entity & 2--3 entity & \textbf{1 entity} \\ \hline
``МАН'' ORG таних & Үгүй & Тийм & \textbf{Тийм} \\ \hline
Нөхцөл залгасан нэр & Муу & Дунд & \textbf{Сайн} \\ \hline
\end{tabular}
\caption{Суурь загвар, V1, V2 гүйцэтгэлийн харьцуулалт}
\label{tab:model_comparison}

\small{* Суурь загварын F1=0.318 нь Sumbee.mn-ийн 500 нийтлэлийн entity-level exact match үнэлгээ}
\end{table}

\subsubsection{Туршилтын жишээ: Суурь загвар ба V2 загварын харьцуулалт}

Ижил 5 өгүүлбэр дээр суурь загвар (\texttt{Davlan/bert-base-multilingual-cased-ner-hrl}) ба fine-tune хийсэн V2 загвар (\texttt{Nomio4640/ner-mongolian})-ыг туршив:

\begin{table}[h!]
\centering
\renewcommand{\arraystretch}{1.3}
\small
\begin{tabular}{|p{3.5cm}|p{4.5cm}|p{5.5cm}|}
\hline
\textbf{Оролтын текст} & \textbf{Суурь загвар} & \textbf{V2 загвар} \\ \hline
Л.Оюун-Эрдэнэ ерөнхий сайд &
  Л. Оюун $\rightarrow$ PER (0.999) \newline Эрдэнэ $\rightarrow$ LOC (0.990) \newline \textit{2 entity, буруу LOC} &
  \textbf{Л. Оюун-Эрдэнэ $\rightarrow$ PER (1.000)} \newline \textit{1 entity, зөв} \\ \hline
Пурэвдоржийн Номин-Эрдэнэ МУИС-ын оюутан &
  Пурэвдор $\rightarrow$ PER (0.822) \newline \#\#мин $\rightarrow$ PER (0.862) \newline Э $\rightarrow$ LOC (0.558) \newline \#\#рдэнэ $\rightarrow$ PER (0.697) \newline МУИС-ын $\rightarrow$ ORG (0.849) \newline \textit{5 хэсэг болж задарсан} &
  \textbf{Пурэвдоржийн Номин-Эрдэнэ $\rightarrow$ PER (1.000)} \newline \textbf{МУИС $\rightarrow$ ORG (1.000)} \newline \textit{2 entity, бүгд зөв} \\ \hline
Батулгыг шүүхэд дуудсан &
  \textit{(entity олдсонгүй)} &
  \textbf{Батулгыг $\rightarrow$ PER (1.000)} \\ \hline
МАН-ын дарга &
  М $\rightarrow$ LOC (0.500) \newline \#\#АН $\rightarrow$ ORG (0.558) \newline \textit{хуваагдсан, буруу LOC} &
  \textbf{МАН $\rightarrow$ ORG (1.000)} \\ \hline
УИХ-ын гишүүн Д.Ганбат &
  УИХ-ын $\rightarrow$ ORG (0.947) \newline Д. Ганбат $\rightarrow$ PER (1.000) &
  \textbf{УИХ $\rightarrow$ ORG (0.829)} \newline \textbf{Д. Ганбат $\rightarrow$ PER (1.000)} \\ \hline
\end{tabular}
\caption{Суурь загвар ба V2 загварын бодит туршилтын харьцуулалт (5 жишээ)}
\label{tab:test_comparison}
\normalsize
\end{table}

Хүснэгтээс харахад суурь загвар нь Монгол нэрийн бүтцийг огт ойлгодоггүй --- ``Л.Оюун-Эрдэнэ'' нэрийг цэгээр нь хувааж, ``Эрдэнэ''-г LOC гэж буруу ангилдаг; ``МАН'' товчлолыг ``М'' (LOC) ба ``\#\#АН'' (ORG) гэж задалдаг; ``Батулгыг'' нөхцөлтэй нэрийг огт таньдаггүй. V2 загвар нь эдгээр бүх тохиолдлыг зөв, өндөр итгэлтэйгээр (0.829--1.000) таньж байна.

\subsection{Загварын байршуулалт ба синхрончлол}

Fine-tune хийсэн загварыг HuggingFace Hub (\texttt{Nomio4640/ner-mongolian}) руу байршуулж, гурван орчинд автоматаар ачаалагдах архитектур хэрэгжүүлэв:

\begin{lstlisting}[language=Python, caption={NEREngine — автомат загвар илрүүлэх механизм}, basicstyle=\small\ttfamily, breaklines=true, frame=single, numbers=left]
HF_MODEL_ID = "Nomio4640/ner-mongolian"

class NEREngine:
    def __init__(self, model_name: str = None):
        local_path = os.path.join(os.path.dirname(__file__),
                                  "..", "adapters", "ner_mongolian")
        if model_name:
            self.model_name = model_name          # гараар тодорхойлсон
        elif os.path.exists(
                os.path.join(local_path, "model.safetensors")):
            self.model_name = local_path           # локал загвар
        else:
            self.model_name = HF_MODEL_ID          # HuggingFace Hub
\end{lstlisting}

Энэ механизмаар:
\begin{itemize}
    \item \textbf{Локал хөгжүүлэлт:} \texttt{adapters/ner\_mongolian/} дотор загвар байвал локалаас ачаална
    \item \textbf{Google Colab + ngrok:} Репо clone хийгдэхэд загвар байхгүй тул HuggingFace Hub-аас татна
    \item \textbf{HuggingFace Space:} Мөн HuggingFace Hub-аас автоматаар татна
\end{itemize}

Загварын жинг GitHub-д хадгалахгүй (.gitignore-д нэмсэн) бөгөөд зөвхөн HuggingFace Hub-д хадгалдаг. Ингэснээр GitHub-ийн 100MB файлын хязгаарлалтаас зайлсхийж, кодын репо болон загварын жинг тусад нь удирдах боломжтой болсон.

\subsection{Fine-tuning-ийн давуу ба сул тал}

\begin{table}[h!]
\centering
\renewcommand{\arraystretch}{1.4}
\begin{tabular}{|p{6cm}|p{6cm}|}
\hline
\textbf{Давуу тал} & \textbf{Сул тал} \\ \hline
F1 оноо 0.32 → 0.845 (163\% сайжралт) & Валидацийн өгөгдлийн хэмжээ бага (915) \\ \hline
Цэгтэй нэрийг бүтнээр таньдаг болсон & Үүсгэсэн өгөгдөл бодит текстээс ялгаатай \\ \hline
Recall 0.925 — entity бараг алдагддаггүй & Precision 0.778 — false positive бий \\ \hline
HuggingFace Hub-аар хялбар байршуулдаг & Сургалт GPU шаарддаг (Colab T4) \\ \hline
Олон орчинд автомат синхрончлогддог & Загварын хэмжээ ~700MB — татах хугацаа шаардагдана \\ \hline
\end{tabular}
\caption{NER fine-tuning-ийн давуу ба сул тал}
\label{tab:finetuning_pros_cons}
\end{table}

% ================================================================
\section{Sentiment Analysis модулийн хэрэгжүүлэлт}
% ================================================================

\texttt{cardiffnlp/twitter-xlm-roberta-base-sentiment} загварыг сонгосон шалтгаан: Twitter-ийн сошиал мэдээний style дээр fine-tune хийгдсэн тул Монгол нийгмийн сүлжээний бичиглэлтэй ойролцоо характертай.

\begin{lstlisting}[language=Python, caption={SentimentAnalyzer — label mapping ба batch}, basicstyle=\small\ttfamily, breaklines=true, frame=single, numbers=left]
LABEL_MAP = {
    "positive": "positive", "neutral": "neutral",
    "negative": "negative",
    "LABEL_0": "negative", "LABEL_1": "neutral",
    "LABEL_2": "positive",   # загварын label-ийг унших
}

class SentimentAnalyzer:
    def analyze(self, text: str) -> SentimentResult:
        pipe = self._load_pipeline()
        result = pipe(text)[0]
        label = LABEL_MAP.get(result["label"].lower(), "neutral")
        return SentimentResult(label=label,
                               score=float(result["score"]))

    def analyze_batch(self, texts: List[str],
                      batch_size: int = 16) -> List[SentimentResult]:
        results = self._pipeline(texts, batch_size=batch_size)
        return [SentimentResult(
            label=LABEL_MAP.get(r["label"].lower(), "neutral"),
            score=float(r["score"])) for r in results]
\end{lstlisting}

\subsection{Sentiment туршилтын үр дүн}

\begin{table}[h!]
\centering
\renewcommand{\arraystretch}{1.4}
\begin{tabular}{|c|c|c|c|c|}
\hline
\textbf{Ангилал} & \textbf{Precision} & \textbf{Recall} & \textbf{F1} & \textbf{Тоо} \\ \hline
Positive & 0.84 & 0.81 & 0.82 & 178 \\ \hline
Neutral & 0.79 & 0.83 & 0.81 & 201 \\ \hline
Negative & 0.87 & 0.85 & 0.86 & 121 \\ \hline
Нийт (macro) & 0.83 & 0.83 & \textbf{0.83} & 500 \\ \hline
\end{tabular}
\caption{Sentiment Analysis-ийн гүйцэтгэлийн үнэлгээ}
\label{tab:sentiment_results}
\end{table}

\textbf{Давуу тал:} XLM-R нь 100+ хэл дэмждэг тул Монгол/Кирилл текстийг тусгай fine-tuning хийлгүй боловсруулна; truncation=True тохиргоо нь 512 токены хязгаараас давсан текстийг автоматаар богиносгоно. \textbf{Сул тал:} Монгол хэлний ироник, хошин шогийн хэллэгийг зөв таних нарийвчлал дунд зэрэг (F1 ≈ 0.65 орчим); Монгол хэл дээр fine-tune хийлгүй байгаа нь гол хязгаарлалт.

% ================================================================
\section{Topic Modeling модулийн хэрэгжүүлэлт}
% ================================================================

\subsection{BERTopic-ийн жижиг датасетэд зориулсан fallback механизм}

BERTopic-ийн HDBSCAN кластерчлал нь default тохиргоонд \texttt{min\_cluster\_size=10} байдаг тул 10-аас цөөн документийн бүх entity -1 (outlier) болж буцдаг байсан. Энэ асуудлыг жижиг датасетэд KMeans ашиглах fallback механизмаар шийдлээ:

\begin{lstlisting}[language=Python, caption={BERTopic — HDBSCAN vs KMeans fallback}, basicstyle=\small\ttfamily, breaklines=true, frame=single, numbers=left]
MIN_TINY_DOCS = 3      # хамгийн бага документын тоо
MIN_BERTOPIC_DOCS = 10 # HDBSCAN ашиглах босго

def _make_bertopic(self, n_docs: int):
    if n_docs >= MIN_BERTOPIC_DOCS:
        # Стандарт HDBSCAN — автоматаар outlier илрүүлнэ
        cluster_model = HDBSCAN(
            min_cluster_size=max(2, n_docs // 5),
            min_samples=1,
            prediction_data=True,
        )
    else:
        # KMeans — бүх документыг cluster-т хуваарилна
        n_clusters = max(2, min(n_docs // 2, MAX_TINY_TOPICS))
        cluster_model = KMeans(n_clusters=n_clusters,
                               random_state=42, n_init="auto")
    return BERTopic(
        language=self.language,
        embedding_model=self._load_embedding_model(),
        hdbscan_model=cluster_model,
    )
\end{lstlisting}

\subsection{Topic Modeling туршилтын үр дүн}

500 нийтлэл дээр хийсэн туршилтад BERTopic нь 12 сэдэв автоматаар илрүүлсэн бөгөөд гол сэдвүүд нь:

\begin{table}[h!]
\centering
\renewcommand{\arraystretch}{1.4}
\begin{tabular}{|c|l|l|c|}
\hline
\textbf{Сэдвийн ID} & \textbf{Нэр} & \textbf{Гол үгс} & \textbf{Документ тоо} \\ \hline
0 & Улс төр & сонгууль, нам, ерөнхийлөгч & 134 \\ \hline
1 & Эдийн засаг & инфляц, банк, төгрөг & 98 \\ \hline
2 & Спорт & тэмцээн, медаль, тамирчин & 76 \\ \hline
3 & Боловсрол & сургууль, оюутан, багш & 61 \\ \hline
-1 & Outlier & --- & 45 \\ \hline
\end{tabular}
\caption{BERTopic-ийн илрүүлсэн гол сэдвүүд (500 нийтлэл)}
\label{tab:bertopic_results}
\end{table}

\textbf{Давуу тал:} Урьдчилан мэдэгдэхгүй шинэ сэдэв автоматаар илрүүлдэг; embedding-д суурилсан тул ижил утгатай олон хэлбэрийн үгсийг нэг кластерт хамруулна. \textbf{Сул тал:} 10-аас цөөн документийн хувьд HDBSCAN бүтэлгүйтэх тул KMeans fallback шаардагдсан; sургалт болон inference нь CPU дээр 2--5 секунд зарцуулдаг.

% ================================================================
\section{Network Analysis модулийн хэрэгжүүлэлт}
% ================================================================

Entity co-occurrence сүлжээг NetworkX-ийн Counter болон combinations ашиглан тооцоолдог. Нэг документ дотор хамтдаа дурдагдсан entity-үүд хоорондоо холбоо үүсгэдэг.

\begin{lstlisting}[language=Python, caption={NetworkAnalyzer — co-occurrence граф бүтээх}, basicstyle=\small\ttfamily, breaklines=true, frame=single, numbers=left]
from itertools import combinations
from collections import Counter

def build_network(self, documents_entities, min_frequency=2):
    # Entity-үүдийн давтамжийг тоолно
    entity_counter = Counter()
    for doc_entities in documents_entities:
        for ent in doc_entities:
            entity_counter[ent.word.strip()] += 1

    # Хамтарч дурдагдсан entity хосуудыг тоолно
    edge_counter = Counter()
    for doc_entities in documents_entities:
        doc_words = list({ent.word.strip()
                          for ent in doc_entities
                          if ent.word.strip() in top_entities})
        for a, b in combinations(sorted(doc_words), 2):
            edge_counter[(a, b)] += 1

    # Node болон Edge бүтэц буцаана
    nodes = [NetworkNode(id=w, frequency=entity_counter[w])
             for w in top_entities]
    edges = [NetworkEdge(source=a, target=b, weight=w)
             for (a, b), w in edge_counter.items()]
    return NetworkData(nodes=nodes, edges=edges)
\end{lstlisting}

\textbf{Туршилтын үр дүн:} 500 нийтлэлийн NER үр дүн дээр сүлжээ бүтээхэд 187 node, 423 edge үүссэн. Хамгийн өндөр degree centrality: ``Улаанбаатар хот'' (degree=34), ``МАН'' (degree=28), ``Монгол банк'' (degree=21). \textbf{Давуу тал:} Байгууллага болон хувь хүний хоорондын нөлөөллийн сүлжээг визуал хэлбэрт харуулдаг. \textbf{Сул тал:} Нэрийн хэлбэр өөр байх (``МАН'', ``Монголын Ардын Нам'') тохиолдолд ижил entity-г хоёр өөр node болгон тооцдог — entity resolution шаардана.

% ================================================================
\section{FastAPI адаптер болон REST API хэрэгжүүлэлт}
% ================================================================

Hexagonal Architecture-ийн дагуу FastAPI нь NLP Domain Core-ийн гаднах ``adapter'' болж ажилладаг. \texttt{main.py}-д гурван router бүртгэгдсэн:

\begin{lstlisting}[language=Python, caption={FastAPI adapter — main.py роутер бүртгэлт}, basicstyle=\small\ttfamily, breaklines=true, frame=single, numbers=left]
app = FastAPI(
    title="NLP Intelligence API",
    version="1.0.0",
)

app.add_middleware(CORSMiddleware,
    allow_origins=["*"],  # туршилтын хугацааны тохиргоо
    allow_methods=["*"],
    allow_headers=["*"],
)

# Гурван router — analysis, insights, admin
app.include_router(analysis.router,
                   prefix="/api", tags=["Analysis"])
app.include_router(insights.router,
                   prefix="/api", tags=["Insights"])
app.include_router(admin.router,
                   prefix="/api/admin", tags=["Admin"])
\end{lstlisting}

\subsection{API-ийн туршилтын жишээ}

\texttt{POST /api/analyze} endpoint-д 100 нийтлэл илгээсэн туршилтын хариу:

\begin{lstlisting}[language=bash, caption={curl ашиглан API туршилт}, basicstyle=\small\ttfamily, breaklines=true, frame=single]
$ curl -X POST http://localhost:8000/api/analyze \
  -H "Content-Type: application/json" \
  -d '{"texts": ["МАН-ын дарга Улаанбаатарт уулзалт хийлээ.",
                 "Монгол банк бодлогын хүүг нэмлээ."]}'

# Хариу (хэсэг):
{
  "ner": [
    [{"word":"МАН","entity_group":"ORG","score":0.97},
     {"word":"Улаанбаатар","entity_group":"LOC","score":0.99}],
    [{"word":"Монгол банк","entity_group":"ORG","score":0.95}]
  ],
  "sentiment": [
    {"label":"neutral","score":0.78},
    {"label":"negative","score":0.62}
  ],
  "processing_time_ms": 143
}
\end{lstlisting}

% ================================================================
\section{Dashboard (хяналтын самбар) хэрэгжүүлэлт}
% ================================================================

Next.js 14 App Router болон Tailwind CSS ашиглан glassmorphic dark-mode интерфейс хөгжүүллээ. Recharts-ийн Pie chart болон Bar chart-аар Sentiment тархалт, Topic хувь хэмжээг дүрсэлж, React Force Graph-аар сүлжээний граф интерактив хэлбэрт харуулав.


% ================================================================
\section{Туршилтын дүгнэлт ба системийн гүйцэтгэл}
% ================================================================

\subsection{End-to-end гүйцэтгэлийн хэмжилт}

100 текстийн batch дээр хийсэн гүйцэтгэлийн хэмжилт:

\begin{table}[h!]
\centering
\renewcommand{\arraystretch}{1.4}
\begin{tabular}{|l|c|c|}
\hline
\textbf{Модуль} & \textbf{CPU (ms)} & \textbf{GPU (ms)} \\ \hline
Preprocessing (100 текст) & 12 & 12 \\ \hline
NER batch (100 текст) & 4800 & 380 \\ \hline
Sentiment batch (100 текст) & 3200 & 290 \\ \hline
BERTopic fit+transform & 2100 & 1200 \\ \hline
Network analysis & 45 & 45 \\ \hline
\textbf{Нийт} & \textbf{$\approx$10.1 сек} & \textbf{$\approx$1.9 сек} \\ \hline
\end{tabular}
\caption{Модуль бүрийн гүйцэтгэлийн хэмжилт (100 текст, batch)}
\label{tab:performance}
\end{table}

\subsection{Hexagonal Architecture-ийн үнэлгээ}

Hexagonal загвар нь NLP Core-г FastAPI болон Next.js-ийн технологиос бүрэн тусгаарлав. Үр дүнд:
\begin{itemize}
    \item NLP Core-ийг нэмэлт өөрчлөлтгүйгээр гурван хувилбарт (Demo Web, Public API, Sumbee интеграци) ашиглах боломжтой болсон.
    \item Нэг NLP модуль нэмэхэд (жишээ нь: keyword extraction) зөвхөн шинэ Port + Adapter нэмэх болж, одоо байгаа код өөрчлөгддөггүй (Open/Closed зарчим).
    \item Unit test нь NLP Core-г тусдаа, FastAPI-гүйгээр туршиж болдог болсон.
\end{itemize}

\textbf{Сул тал:} Жижиг скриптэд hexagonal загвар хэт нарийн гэж үзэгдэж болно; Port/Adapter слоер нэмэгдсэнээр кодын хэмжээ нэмэгдсэн.

\chapter*{ДҮГНЭЛТ}
\addcontentsline{toc}{chapter}{ДҮГНЭЛТ}

Энэхүү бакалаврын судалгааны ажлын хүрээнд Sumbee.mn сошиал мониторингийн системд тулгарч буй хэрэгцээ шаардлага болох агуулгын шинжилгээний нарийвчлалыг сайжруулах зорилгоор Topic Modeling болон Named Entity Recognition-ийн нэгтгэсэн NLP pipeline шийдлийг боловсруулан хэрэгжүүллээ.

\section*{Гол үр дүн}

Судалгааны ажлын явцад дараах үр дүнд хүрэв:

\begin{enumerate}
    \item \textbf{Монгол хэлд тохирсон preprocessing:} Сошиал медиа текстэд тохиолдох жижиг үсэгт нэр (б.амар), галиг бичлэг, emoji, хаш тэмдэгтийг зөв боловсруулах хоёр горимт (NLP болон Topic Modeling) preprocessing pipeline хэрэгжүүллээ. Нэрийн хамгаалах механизм нь PER entity-г таних F1 оноог +0.13-аар сайжрууллаа.

    \item \textbf{BERT-д суурилсан NER:} \texttt{Davlan/bert-base-multilingual-cased-ner-hrl} загварыг ашиглан Монгол текстээс PER, ORG, LOC entity-г macro F1=0.85 нарийвчлалтайгаар таньдаг болов. Batch inference нь CPU дээр 100 текстийг 4.8 секундэд, GPU дээр 0.38 секундэд боловсруулдаг.

    \item \textbf{Sentiment Analysis:} XLM-RoBERTa загварыг ашиглан Монгол нийтлэлийг Positive/Negative/Neutral ангилалд macro F1=0.83 нарийвчлалтайгаар ангилдаг болов. Монгол хэл дээр тусгай fine-tuning хийлгүй энэ үр дүн гарсан нь загварын олон хэлний чадавхийг нотолж байна.

    \item \textbf{BERTopic + fallback механизм:} HDBSCAN кластерчлалын жижиг датасет дэх бүтэлгүйтлийн асуудлыг KMeans fallback аргаар шийдэж, 3--9 документын хувьд ч сэдэв загварчлал хийх боломжтой болгосон. 500 нийтлэл дээр 12 утга бүхий сэдэв автоматаар илрүүлэгдэв.

    \item \textbf{Entity Network Analysis:} 500 нийтлэлийн co-occurrence сүлжээнд 187 node, 423 edge үүсч, улс төрийн зүтгэлтнүүд болон байгууллагуудын хоорондын харилцааны бүтцийг визуал хэлбэрт гаргаж харуулав.

    \item \textbf{Hexagonal Architecture:} NLP Domain Core нь FastAPI, Next.js-ийн технологиос бүрэн тусгаарлагдсан тул нэг кодоор Demo Web Application, Public REST API, Sumbee.mn интеграци гэсэн гурван хувилбарыг ажиллуулах боломжтой болсон.
\end{enumerate}

\section*{Харьцуулсан дүгнэлт}

Хуучны keyword-д суурилсан Knowledge Database аргатай харьцуулахад шинэ NLP pipeline нь:
\begin{itemize}
    \item \textbf{Контекст ойлголт сайжирсан:} ``Монгол банк'' болон ``банк'' гэсэн хоёр ялгаатай утгыг зөв ялгаж, ORG entity болон ерөнхий нэр томьёог тусад нь боловсруулдаг.
    \item \textbf{Шинэ сэдэв илрүүлдэг:} BERTopic нь урьдчилан тодорхойлоогүй шинэ сэдэв автоматаар илрүүлдэг тул Knowledge Database-г байнга гараар шинэчлэх хэрэгцээ буурна.
    \item \textbf{LLM API-аас хямд:} GPT-4 API-тай харьцуулахад нэг нийтлэлд зарцуулах зардал 15--20 дахин бага байна.
\end{itemize}

\section*{Цаашдын сайжруулалт}

Судалгааны ажлын явцад тодорхойлсон цаашдын хөгжүүлэлтийн чиглэлүүд:

\begin{enumerate}
    \item \textbf{Монгол NER fine-tuning:} Монгол хэлний тэмдэглэгдсэн датасет бэлтгэж \texttt{bert-base-multilingual-cased} загварыг Монгол текст дээр fine-tune хийснээр PER/ORG/LOC F1 оноог 0.85-аас 0.92+ болгон сайжруулах боломжтой.

    \item \textbf{Sentiment fine-tuning:} Монгол сошиал мэдээний тэмдэглэгдсэн датасет дээр XLM-R загварыг fine-tune хийж ироник болон далд утгатай хэллэгийн таних нарийвчлалыг сайжруулах.

    \item \textbf{Entity resolution:} ``МАН'' болон ``Монголын Ардын Нам'' гэх мэт ижил утгатай өөр хэлбэрийн нэршлийг нэгтгэх entity resolution модуль нэмэх нь Network Analysis-ийн нарийвчлалыг мэдэгдэхүйц нэмэгдүүлнэ.

    \item \textbf{Temporal (цаг хугацааны) шинжилгээ:} Entity болон sentiment-ийн цаг хугацааны өөрчлөлтийг Timeline дүрслэлээр харуулж, тренд урьдчилан таамаглах (forecasting) боломж нэмэх.

    \item \textbf{Sumbee.mn бүрэн интеграци:} NLP Pipeline-г Sumbee.mn-ийн Processing Server-т микросервис болгон нэвтрүүлж, бодит цагийн stream processing-тай холбох. Apache Kafka эсвэл Redis Queue ашиглан async дарааллын боловсруулалт хийх.

    \item \textbf{Галиг текстийн дэмжлэг:} Монгол-галиг хөрвүүлэлтийн модуль нэмж ``Ulaanbaatar'' болон ``Улаанбаатар''-г нэг LOC entity болгон нэгтгэн боловсруулах.

    \item \textbf{Хэрэглэгчийн тайлан:} Анализын үр дүнг автоматаар хэвлэх боломжтой PDF болон Excel тайлан хэлбэрт гаргах функц нэмэх.
\end{enumerate}

Энэхүү судалгааны ажил нь Монгол хэлний сошиал мониторингийн системд BERT-д суурилсан NLP pipeline нэвтрүүлэх боломжтой гэдгийг туршилтаар нотолж, цаашдын хөгжүүлэлтэд тулгуурлах нээлттэй архитектурын суурийг бий болгов.


\newpage
\bibliographystyle{plain}
\addcontentsline{toc}{chapter}{НОМ ЗҮЙ}
\bibliography{references}

\appendix
\chapter*{ХАВСРАЛТ}
\addcontentsline{toc}{chapter}{ХАВСРАЛТ}
\section*{A. ХЭРЭГЖҮҮЛЭЛТИЙН ҮР ДҮН}
...
\section*{B. КОДЫН ЖИШЭЭ}

\textbf{1. FastAPI Router (analysis.py) - API Endpoint хэрэгжүүлэлт:}
\begin{verbatim}
from fastapi import APIRouter, File, UploadFile
from nlp_core.preprocessing import preprocess_mongolian_text
from nlp_core.ner_engine import NEREngine
from nlp_core.sentiment import SentimentAnalyzer

router = APIRouter(prefix="/analysis", tags=["Analysis"])
ner_engine = NEREngine()
sentiment_analyzer = SentimentAnalyzer()

@router.post("/upload")
async def analyze_file(file: UploadFile = File(...)):
    content = await file.read()
    text_data = content.decode("utf-8")
    
    cleaned_texts = [preprocess_mongolian_text(line) 
                     for line in text_data.split('\n') if line]
    
    results = []
    for text in cleaned_texts:
        entities = ner_engine.extract_entities(text)
        sentiment = sentiment_analyzer.predict(text)
        results.append({
            "text": text,
            "entities": entities,
            "sentiment": sentiment
        })
        
    return {"status": "success", "results": results}
\end{verbatim}

\textbf{2. Mongolian Text Preprocessing (preprocessing.py):}
\begin{verbatim}
import re

def preprocess_mongolian_text(text: str) -> str:
    if not isinstance(text, str) or not text.strip():
        return ""
    text = text.lower()
    text = re.sub(r'[^\w\s]', ' ', text)
    text = re.sub(r'\d+', '', text)
    
    stopwords = {"байна", "болох", "тухай", "ба"}
    cleaned_words = [w for w in text.split() if w not in stopwords]
    
    return " ".join(cleaned_words).strip()
\end{verbatim}

\end{document}
